\documentclass[11pt,letterpaper]{article}
\usepackage[T1]{fontenc}
\usepackage[utf8]{inputenc}
\usepackage{lmodern,microtype}
\usepackage[margin=0.91in,headheight=15pt]{geometry}
\usepackage{amsmath,amssymb,amsthm,mathtools}
\usepackage{booktabs,longtable,array,tabularx,enumitem}
\usepackage{xcolor,fancyhdr,listings,xurl,needspace}
\usepackage[colorlinks=true,linkcolor=blue!35!black,citecolor=blue!35!black,urlcolor=blue!45!black]{hyperref}
\usepackage{bookmark}
\definecolor{ink}{RGB}{25,47,69}
\definecolor{pale}{RGB}{245,247,249}
\setlength{\emergencystretch}{3em}
\setlength{\parskip}{0.3em}
\setlist{itemsep=3pt,topsep=5pt,leftmargin=1.6em}
\newtheorem{theorem}{Theorem}[section]
\newtheorem{lemma}[theorem]{Lemma}
\newtheorem{proposition}[theorem]{Proposition}
\newtheorem{corollary}[theorem]{Corollary}
\theoremstyle{definition}\newtheorem{definition}[theorem]{Definition}
\theoremstyle{remark}\newtheorem{remark}[theorem]{Remark}
\newcommand{\ipra}{\mathsf{iPRA}}
\newcommand{\D}{\mathsf D}
\newcommand{\Q}{\mathsf Q}
\newcommand{\Eq}{\mathsf{Eq}}
\newcommand{\Tr}{\mathsf{Tr}}
\newcommand{\FV}{\operatorname{FV}}
\newcommand{\llbracket}{\mathopen{[\![}}
\newcommand{\rrbracket}{\mathclose{]\!]}}
\newcommand{\file}[1]{\expandafter\nolinkurl\expandafter{\detokenize{#1}}}
\newcommand{\Acal}{\mathcal A}
\newcommand{\Ccal}{\mathcal C}
\lstset{basicstyle=\ttfamily\footnotesize,columns=fullflexible,keepspaces=true,
 breaklines=true,showstringspaces=false,frame=single,rulecolor=\color{black!15},
 backgroundcolor=\color{pale},aboveskip=8pt,belowskip=8pt}
\pagestyle{fancy}\fancyhf{}
\fancyhead[L]{\small A certified zero removes a computation opcode}
\fancyhead[R]{\small 10 November 2026}
\fancyfoot[C]{\thepage}
\hypersetup{pdftitle={A Certified Zero Removes a Computation Opcode: Smaller Complete Presentations of iPRA},pdfauthor={Research prepared for Vladimir Reshetnikov}}
\title{\color{ink}\bfseries A Certified Zero Removes\\a Computation Opcode\\[8pt]
\large Smaller complete presentations of intuitionistic\\primitive recursive arithmetic}
\author{Research prepared for Vladimir Reshetnikov}
\date{10 November 2026}
\begin{document}
\maketitle
\begin{abstract}
A scalar-iteration proof checker can use the initializer itself to
distinguish computation from paired introduction. This eliminates a
separate computation tag, provided direct computation is restricted to
initializers with nonempty syntax. The restriction loses no theorem:
the fixed term $C=I(2,1;v.0)$ has a direct permitted proof of $C=0$.
Replacing a missing zero initializer, a missing zero predecessor, or both
by $C$, and transporting this equality once, recovers every excluded
computation equation. Each excluded computation leaf is replaced by
three nodes. The numerical language, native formulas, and every native
sequent are preserved exactly.

The inherited scalar-iteration and paired-binder presentation is stated
explicitly, together with its arithmetic and logical reconstruction and
both interpretation composites with full first-order iPRA. The new
restriction therefore preserves whole-domain bi-interpretability and
equiconsistency. Complete implementations, proof producers, positive
and adversarial certificates, source inventories, and offline reproduction
scripts accompany the report. Gzip level-nine sizes are reported only as
auxiliary measurements.

The comparison covers all 13 current draft ZIPs at commit
\file{fb2dbd2a33c6b132671a54b029d0bf3cd98c633d}, including 370 distinct
Python sources. Both the supplied frontier and stronger matched-driver
controls are measured. One runtime-qualified implementation is below
all three matched coordinate minima; the recommended explicit-list
implementation has smaller source bytes and avoids the malformed-string
qualification. The improvement is a proved restriction and a smaller
checker, not an absolute minimality theorem or the requested major new
arithmetic foundation.
\end{abstract}
\begin{center}
\fcolorbox{ink}{pale}{\begin{minipage}{0.92\linewidth}
\textbf{What is established.} A new admissible computation restriction,
an eliminated packet opcode, complete smaller implementations, and
identity preservation of the existing full-iPRA interpretation.
\textbf{What remains open.} An absolute simplest axiomatization and a
significant foundational breakthrough. The mathematical proofs are
written proofs, not Lean or Rocq formalizations.
\end{minipage}}
\end{center}
\clearpage
\setcounter{tocdepth}{1}
\hypersetup{bookmarksdepth=2}
\begingroup\small\setlength{\parskip}{0pt}\tableofcontents\endgroup
\clearpage

\section{Result, comparison, and attribution}
\label{sec:result}

The main change is small enough to explain without the encoding.
Computation packets already need an initializer. Ensure that its syntax
is nonempty, and its presence can select computation; paired-introduction
packets leave that field empty. The difficulty is mathematical: zero is
an indispensable initializer. The closed identity
\[
                  \boxed{I(S(S0),S0;v.0)=0}
\]
solves that difficulty using one computation whose predecessor and
initializer are both nonempty. No earlier zero-initializer theorem is
used. Equality transport then restores the missing cases.

\subsection{Exact whole-program comparison}

The source snapshot is \cite{corpus}. The recursive inventory contains
13 direct archives, 2,701 Python occurrences, and 370 distinct raw source
hashes. All 31 distinct programs at or below at least one supplied
coordinate minimum were read and classified: three are complete
checkers, and 28 are supporting programs. Tiny build scripts and test
drivers that call a checker are not counted as proof checkers.

The supplied coordinate minima are 268 significant normalized tokens,
300 AST nodes, and 616 normalized UTF-8 bytes. Different files attain
them. The often quoted larger headline values in individual reports do
not include every ancillary or matched implementation in this snapshot.
The seven additional matched controls remove only the obsolete
\texttt{print('REJECT'); raise} wrapper from selected affine and arity
programs. These controls establish the stronger comparison envelope
264 tokens, 300 nodes, and 601 bytes. Their exact AST transformation,
original bytes, hashes, and verification receipts are included.


% BEGIN evidence/metrics_table.tex
\begin{center}
\small
\begin{tabular}{lrrrr}
\toprule
Complete implementation & Tokens & AST & Bytes & Gzip-9\\
\midrule
Supplied coordinate minima & 268 & 300 & 616 & ---\\
Matched coordinate minima & 264 & 300 & 601 & ---\\
\midrule
New explicit list test & 264 & 308 & 566 & 351\\
New byte profile & 264 & 310 & 558 & 345\\
New joint profile (qualified) & 262 & 296 & 598 & 350\\
\bottomrule
\end{tabular}
\end{center}
All new raw-source lengths equal their normalized lengths. 
The two envelope rows are coordinatewise minima from different files; 
each new implementation row is one complete file.

% END evidence/metrics_table.tex


\file{code/checker.py} is the recommended implementation: it uses an
explicit \texttt{isinstance(a,list)} test. The byte profile uses the
inherited exact-JSON-type predicate \texttt{a*0 == []}. The joint profile
uses the inherited exception-based traversal; it has the smallest
combined structural counts but a qualification on malformed strings.
These are distinct complete programs. No row combines measurements
from different programs.

The joint profile is strictly below every coordinate of both envelopes.
That statement includes its runtime qualification; it must not be
silently attributed to the recommended list-based checker. The
recommended and byte profiles primarily improve source bytes. The
observed envelopes are achieved upper bounds on an unknown optimum,
not mathematical lower bounds on all theories or implementations.

\subsection{Measurement convention}

The whole module is counted: imports, JSON input, syntax and scope
checks, arithmetic and logical rules, and the acceptance driver.
Normalization removes genuine docstrings and comments through AST
reconstruction, uses \texttt{ast.unparse}, and appends one newline.
Tokens exclude whitespace, comments, indentation markers and end markers.
AST counts include the full module. CPython 3.12.14 and zlib 1.3.2 were
used; another implementation can change normalization or compression.

Both raw and normalized source are compressed by
\begin{lstlisting}[language=Python]
gzip.compress(source_bytes, compresslevel=9, mtime=0)
\end{lstlisting}
including the gzip wrapper and no filename. The exact \texttt{.py.gz}
streams accompany the sources. Gzip is informational; it was not the
optimization objective. Python and its standard JSON implementation
remain the common execution environment used by all compared programs.

\subsection{New and inherited components}

The initialized scalar iterator, Choice-DI logic, arithmetic bootstrap,
uniform predecessor packets, unary/triplicate formula representation,
arity-based scope traversal, and generic Python representation profiles
are inherited from \cite{arity,uniform,constructed,generated,transient}.
The new mathematical step restricts direct initializers and proves the
restriction admissible with an independently constructed zero. The new
certificate protocol uses that data field in place of an arithmetic tag.
The last source reduction shares the already required assumption
assignment across its three cases; it introduces no new rule.

An explicit ablation isolates the contribution of initializer dispatch.
Restoring only the old \texttt{m is None} discriminator in the final
three programs gives respectively 574, 566 and 606 bytes. The new
discriminator saves exactly eight bytes, two tokens and three AST nodes
in each corresponding program. These newly constructed ablations are
separate from the seven driver-only controls and are included in
\file{reference/ablations/}. Thus the larger total gain against the
supplied corpus is not attributed entirely to the normalization theorem.

A repository-relative audit found no initializer-based computation
dispatch among 370 distinct Python sources and 386 distinct textual
sources. Earlier drafts do use proved nonempty representations of zero
for witnesses and predecessors. That general technique is credited;
its application to initializer dispatch and the simultaneous repair here
are the new step. No priority claim outside this corpus is made.

\section{The complete native theory}
\label{sec:inventory}

Here iPRA means full first-order intuitionistic arithmetic in a language
of fixed finite primitive-recursive descriptions, with their defining
equations and quantifier-free induction. Atomic induction is a customary
equivalent formulation in the full primitive-recursive language; see
Salehi's definition \cite{salehi}. The internal conversion needed for our
smaller language is proved in the attributed appendices.

\subsection{Numerical expressions and restricted computation}

There is one numerical sort, with finite terms
\[
             t ::= x\mid 0\mid St\mid I(n,a;v.h).
\]
Only $v$ in $h$ is bound. All other arguments of $h$ are passive
parameters, fixed during iteration. Each fixed previously constructed
update description gives a first-order function symbol; there is no
quantification over functions and no uniformly total evaluator of all
program codes.

There are exactly two nonlogical computation patterns:
\begin{align}
 I(0,a;v.h)&=a &&\text{if $a$ is not the literal zero term},\tag{R0}\label{eq:r0}\\
 I(Sn,a;v.h)&=h[I(n,a;v.h)/v]
 &&\text{if neither $n$ nor $a$ is literal zero}.\tag{RS}\label{eq:rs}
\end{align}
The side conditions are metasyntactic. A permitted term may denote zero.
Substitution in the update is capture avoiding. The restriction makes
arbitrary numerical substitution an admissible transformation rather
than a closure property of the displayed primitive axiom list: a
substitution producing a forbidden literal zero is repaired by the
next section's construction.

\subsection{Object formulas and auxiliary closures}

Object formulas are
\[
              A ::= (s=t)\mid\D_{x,y}A.
\]
An auxiliary closure $\Q_yB$ may occur in contexts and certificates.
It may also be admitted as a definable formula constructor, since it is
interderivable with $\D_{u,y}B$ for unused $u$. Ordinary readings are
\[
 (\Q_yB)^\circ=B^\circ(0)\to\forall y\,B^\circ(y),\qquad
 (\D_{x,y}B)^\circ=
 \exists x\bigl(B^\circ(x,0)\to\forall y\,B^\circ(x,y)\bigr).
\]
Contexts are finite reusable lists. The designated native contradiction
is $0=S0$. No standard propositional connective or quantifier is a
separate native primitive.

\subsection{All primitive logical rules}

Hypothesis, equality transport, and closure specialization are
\[
 \frac{}{\Gamma\vdash A}\ (A\in\Gamma),\qquad
 \frac{\Gamma\vdash s=t\quad\Gamma\vdash F(s)}
      {\Gamma\vdash F(t)},\qquad
 \frac{\Gamma\vdash\Q_yB(y)\quad\Gamma\vdash B(0)}
      {\Gamma\vdash B(t)}.
 \tag{Hyp, Tr, QE}
\]
$F$ is a fixed finite native formula template; the transport rule
permits transport through binders. Paired elimination is
\[
 \frac{\Gamma\vdash\D_{x,y}B(x,y)\qquad
       \Gamma,\Q_yB(b,y)\vdash A}{\Gamma\vdash A}.
 \tag{DE}
\]
The eigenparameter $b$ is absent from $\Gamma,A$ and the fixed parameters
of the matrix. Paired introduction has the two alternatives
\[
 \frac{\Gamma,B(t,0)\vdash B(t,Sz)}
      {\Gamma\vdash\D_{x,y}B(x,y)},\qquad
 \frac{\Gamma,E(t,z)\vdash E(t,Sz)}
      {\Gamma\vdash\D_{x,y}E(x,y)}.
 \tag{DI-case, DI-progress}
\]
Progress is restricted to a literal equation $E$, not a quantified
formula or an equation beneath a binder. In both alternatives $z$ is
fresh for the retained context, witness and passive parameters, and the
displayed local assumption is discharged.

The computation equations can be axiom leaves. The implementation uses
the equivalent local-assumption form
\[
     \frac{\Gamma,E\vdash A}{\Gamma\vdash A}
     \quad(E\text{ a permitted R0 or RS instance}).\tag{Supply}
\]
A computation leaf is Supply followed by Hyp; Supply is an axiom leaf
followed by admissible cut. These are alternative proof layouts for the
same theory, not extra arithmetic axioms.

Weakening, exchange, contraction, renaming, capture-free numerical
substitution, and ordinary-assumption cut are admissible. They follow
by translating to the unrestricted identical-sequent calculus, applying
the usual finite structural transformation, and repairing computations.
There is no primitive general induction, excluded middle, equality
reflection rule, successor injectivity rule, reflexivity axiom, or
coverage axiom. Their needed intuitionistic or arithmetic instances
are derived as described below.


% BEGIN theory/seed_normalization.tex
% New admissibility argument for initializer-driven computation dispatch.
% Inclusion fragment; requires amsmath and amssymb.

\section{Using the initializer to select computation}
\label{sec:seed-dispatch}

The initialized scalar iterator has the usual equations
\[
 I(0,a;v.h)=a,\qquad
 I(Sn,a;v.h)=h[I(n,a;v.h)/v].
 \tag{I0, IS}
\]
The following restriction removes a computation opcode from the
certificate interface without changing these operations or the native
consequence relation.  It concerns the \emph{syntax} of an initializer,
not whether that initializer denotes a positive natural number.

\paragraph{A fixed nonempty representation of zero.}
Let
\[
               C:=I(S(S0),S0;v.0).
 \tag{Auxiliary zero}\label{eq:dispatch-zero-term}
\]
One successor computation, at predecessor $S0$, proves
\[
                              C=0.\label{eq:dispatch-zero}
 \tag{Zero certificate}
\]
The predecessor and initializer in this computation are both nonempty
canonical terms.  The update is the constant zero, so the right side of
the computation is exactly zero.  This construction does not use a
previous computation with initializer zero, a zero-predecessor repair,
an induction, or a presumed identity for a new numerical constructor.

\paragraph{Restricted computation system.}
In canonical numerical serialization, use scalar $0$ for zero and
nonempty arrays for successor and iteration.  Consider the same native
calculus with its direct computation leaves restricted as follows:

\begin{itemize}
\item I0 may be used directly only when its initializer is not literal
  zero.
\item IS may be used directly only when both its initializer and its
  supplied predecessor are not literal zero.
\end{itemize}

The predecessor restriction is already present in the inherited
predecessor-specified interface.  The new restriction is on the
initializer.  A parameter or a nonempty term is permitted even when its
value happens to be zero.  In particular, $C$ is permitted.  Keep the
ordinary native equality-transport rule and all other logical rules
unchanged.

\paragraph{Admissibility theorem.}
Every unrestricted I0 and IS instance has a finite derivation in this
restricted system.  Consequently the restricted and unrestricted
calculi prove exactly the same native sequents, with exactly the same
numerical and formula languages.

\paragraph{Proof.}
Instances satisfying the stated syntactic conditions are direct.
For a base computation with initializer zero, take the fresh-variable
equation motive
\[
                  E(w):=\bigl(I(0,w;v.h)=w\bigr).
\]
The restricted base equation proves $E(C)$.  Equality transport along
the independently proved $C=0$ gives $E(0)$, which is precisely the
required unrestricted base instance.

For a successor computation with permitted predecessor $n$ but zero
initializer, take
\[
 E(w):=\left(I(Sn,w;v.h)=h[I(n,w;v.h)/v]\right).
\]
The restricted successor equation proves $E(C)$, and the same transport
gives $E(0)$.  The displayed variable $w$ is fresh for $n$, $h$, and all
retained assumptions.  Only the indicated initializer occurrences are
substituted; unrelated zero occurrences or passive arguments in $h$
remain fixed.

For a zero predecessor but a permitted initializer $a$, use instead
\[
 E(w):=\left(I(Sw,a;v.h)=h[I(w,a;v.h)/v]\right).
\]
Again $E(C)$ is a permitted successor computation and transport gives
the desired instance $E(0)$.

If both predecessor and initializer are zero, a single shared motive
repairs them simultaneously:
\[
 E(w):=\left(I(Sw,w;v.h)=h[I(w,w;v.h)/v]\right).
\]
Both supplied values in $E(C)$ are nonempty, and transport along $C=0$
gives exactly $E(0)$.  Bound state variables are renamed before these
substitutions, so no capture occurs.

Each repaired computation leaf therefore uses one permitted computation
for $C=0$, one permitted computation for $E(C)$, and one transport
inference.  Replacing computation leaves in a native derivation proves
the inclusion from the unrestricted calculus to the restricted one;
the reverse inclusion is immediate.  Contexts and final sequents are
unchanged throughout.\hfill$\square$

This is a three-node repair when computation leaves and transport are
counted as one node each and the proof of $C=0$ is copied.  If a tree
derivation has $N$ nodes and $r$ excluded computation leaves, the stated
replacement has exactly $N+2r$ nodes.  If a DAG proof retains one shared
node proving the closed lemma $C=0$, the repaired part instead has
$2r+1$ nodes, giving $N+r+1$ distinct nodes when $r>0$.  A certificate
format with persistent proved assumptions can implement corresponding
sharing by supplying this closed computation once before the affected
continuation.  These are counts for the displayed construction, not
lower bounds on shortest certificates or claims about a particular
producer's serialization.

\paragraph{Why the opcode can disappear.}
Write a four-field packet as $[m,s,u,t]$.  The established elimination
branch uses a truthy major $m$.  In a computation packet, let $s$ be the
predecessor flag/value, $u$ the initializer, and $t$ the update.  A truthy
$u$ selects the computation branch.  The checker constructs the source
counter as $s$ when it is the zero flag and as $Ss$ otherwise, and
generates I0 or IS.  A paired-introduction packet puts zero in the unused
$u$ field.  Thus the dispatch can be written
\begin{verbatim}
    if m:
        ... elimination ...
    elif u:
        ... computation ...
    else:
        ... paired introduction ...
\end{verbatim}
The direct computation branch can no longer receive a literal-zero
initializer.  The admissibility theorem supplies exactly that missing
case.  No test that $u$ denotes a nonzero number is performed or needed.

\paragraph{The fresh-name obligation on noncanonical packets.}
The dispatch change admits computation packets with any false-valued
$m$, including the atomic-progress marker $m=0$.  In the common
continuation, that marker can cause the generated hypothesis to be
opened at the new rigid parameter $b$, rather than at literal zero.
It would be incorrect to assert that the resulting $H(b)$ is uniformly
true.  Use the existing zero-instance invariant instead:
\[
        \vdash H(\vec0),\qquad
        \vec b\cap\operatorname{FV}(\Gamma,A)=\varnothing.
\]
Here $\vec b$ includes the designated new integer name $b$ and every
additional negative fractional name in the newly admitted unit interval
that occurs in the finite raw certificate.  Such names must not be
silently identified with the designated integer.  Given a finite
derivation of $\Gamma,H(\vec b)\vdash A$, rename its later
eigenparameters if necessary, substitute zero for every member of
$\vec b$, and cut the resulting assumption $H(\vec0)$ against its
computation derivation.  This
gives $\Gamma\vdash A$.  The old-floor validation of the retained goal
and context is what establishes the displayed freshness condition.
Consequently the new raw packet possibilities are sound under the same
zero-completion and opening lemmas as the predecessor-specified kernel.
They do not justify silently strengthening a known zero instance into
a uniform computation theorem.

\paragraph{Interpretations and the scope of the improvement.}
The theorem preserves every native sequent through the identity map on
object expressions.  It therefore transfers the inherited whole-domain
bi-interpretation and equiconsistency with iPRA without a new domain,
new equality relation, or new primitive numerical operation.  The
contribution is a removable syntactic restriction and a data-driven
certificate dispatch.  It is an implementation simplification, not a
new arithmetic foundation or a proof of absolute minimality.

Relative to the inspected pinned draft collection, the new step is
using the initializer itself to discriminate computation packets and
proving that the resulting missing computation instances are admissible.
The broader technique of replacing an unavailable zero representation
by a proved nonempty representative is already used in the supplied
drafts, including their predecessor and witness repairs.  The present
argument adapts that attributed technique to eliminate the computation
opcode.  No priority claim outside the inspected corpus is made.

% END theory/seed_normalization.tex


\section{Full iPRA and both interpretation composites}
\label{sec:bridge}

The new normalization preserves every native sequent by the identity
map. It remains necessary to specify the inherited bridge: agreement
on standard numerals or equiconsistency alone would not suffice.

\subsection{Recovering the ordinary logical operations}

Once the omitted computations have been restored, put
\[
 C(v,a,b)=I(v,a;u.b),\quad
 Z(v)=C(v,1,0),\quad p(v)=I(v,0;u.Z(u)),\quad 1=S0.
\]
The use of $C$ here for the conventional three-argument selector is
separate from the closed auxiliary-zero constant of the preceding
section; its arguments distinguish the two notations. Computation gives
$C(0,a,b)=a$, $C(Sv,a,b)=b$, $p(1)=1$, and $p(2)=0$.
For equations define
\[
 G_v(s=t)=(C(v,0,s)=C(v,0,t)),\qquad
 M_v(s=t,r=q)=(C(v,s,r)=C(v,t,q)).
\]
Extend through paired binders by finite syntax recursion, aligning bound
variables and using the proved padding equivalence
$A\dashv\vdash\D_{u,y}G_yA$ when binder depths differ. The full recursive
definition and its proof are in Appendix~\ref{sec:choice-introduction}.
It gives
\[
 \vdash G_0A,\quad G_{Sv}A\dashv\vdash A,\quad
 M_0(A,B)\dashv\vdash A,\quad M_{Sv}(A,B)\dashv\vdash B.
\]
With all new binders fresh and $u$ unused, define
\begin{align*}
 A\Rightarrow_*B&=\D_{u,z}M_z(A,B), &
 A\wedge_*B&=\D_{u,z}G_zM_{p(z)}(A,B),\\
 \exists_*x\,A(x)&=\D_{x,y}G_yA(x), &
 A\vee_*B&=\exists_*x\,M_x(A,B),\\
 \forall_*x\,A(x)&=A(0)\wedge_*\D_{u,z}G_zA(z), &
 \bot_*&=(0=1).
\end{align*}
All ordinary intuitionistic introduction and elimination rules follow
from DI-case, DE, QE and these endpoint identities. The zero conjunct
in the universal definition is necessary. Conjunction projections use
indices one and two. Native explosion first gives every equation by
transport through the selector and then every formula by syntax induction.
These constructions preserve arbitrary retained contexts.

DI-progress with an unused witness, followed by DE and QE, gives
contextual equation induction. DI-case similarly gives zero/successor
cases for arbitrary native formulas. The resulting logic is full
intuitionistic first-order logic; no decidability principle for arbitrary
quantified formulas is introduced.

\subsection{Arithmetic and whole-domain interpretation}

The included arithmetic development has an explicit dependency order:
scalar predecessor; truncated subtraction and a characteristic of actual
equality; quantifier-free induction by characteristic equations; then
pairing and the compilation of arbitrary fixed primitive recursions.
No pairing or equality detector is assumed in the predecessor bootstrap.
All inductions keep passive parameters fixed. The equalities are proved
internally for arbitrary domain elements.

If the target description is
$f(0,\vec a)=g(\vec a)$ and
$f(Sn,\vec a)=h(n,f(n,\vec a),\vec a)$, the proved pairing $\pi$ and
projections $L,H$ give
\begin{align*}
 V(n,\vec a)&=I\bigl(n,\pi(0,g^*(\vec a));
 u.\pi(SL(u),h^*(L(u),H(u),\vec a))\bigr),\\
 f^*(n,\vec a)&=H(V(n,\vec a)).
\end{align*}
Equation induction proves $L(V(n,\vec a))=n$ and the required defining
equations. This is recursion on a fixed finite description in the
metatheory, not a total numerical evaluator of every code.

\begin{theorem}[Bi-interpretability and equiconsistency]
The restricted native calculus presents a theory bi-interpretable and
equiconsistent with full first-order iPRA. Both numerical interpretation
domains are the entire domain, equality is actual equality, and both
domain isomorphisms for the composites have graph $x=y$.
\end{theorem}
\begin{proof}
First recover the unrestricted native consequence relation by
Section~\ref{sec:seed-dispatch}. Interpret each scalar iterator in iPRA
by primitive recursion with an update that ignores its counter. In the
opposite direction use the displayed compiler and the internally proved
arithmetic facts. Ordinary logical rules are recovered by the macros.

For each numerical composite, the translated operation satisfies the
original defining equations. Internal equation induction proves
uniqueness of that fixed recurrence with its parameters held fixed.
Induction on the finite operation description proves equality of each
operation with its double translation, including coherence under
substitution of actual arguments.

The formula composites include every native formula, not only formulas
in the image of the ordinary logical translation. In particular,
\[
 \D_{x,y}B\dashv\vdash
 \exists_*x\bigl(B(x,0)\Rightarrow_*\forall_*yB(x,y)\bigr),
\]
and the analogous closure composite holds. Forward, DE supplies the
fresh witness and its clause; under the zero instance use QE, and then
introduce universal, implication and existential. Backward, eliminate
the existential, retain the zero instance in DI-case, and obtain each
successor instance by implication and universal elimination. External
induction on formulas finishes both composites.

In the ordinary realization, include the explicit falsity bridge
$(0=S0)\to\bot$ as well as numerical coverage; the native derived
explosion and cases supply them. Positive equations alone would not
justify this bridge. Contradiction is preserved in both directions, so
the effective finite proof translations establish equiconsistency.
This is not a claim that either theory proves its own consistency.
\end{proof}


% BEGIN theory/checker.tex
\section{The complete checker and its interpretation}
\label{sec:checker}

The recommended program below is the entire trusted Python source.
No producer, evaluator, theorem database, measurement script or service
is imported. It uses the inherited arity representation \cite{arity}
with the new initializer dispatch and shared assumption assignment.

\begin{lstlisting}[language=Python]
import json
k = 0
i = lambda a, t=None, d=0: [i(b, t, d + 5 - len(a)) for b in a] if isinstance(a, list) else a if k <= a < d else d % 2 and i(t, d=3)

def c(a, g):
    global k
    if i(a) in g:
        return
    m, s, u, t = q()
    if m:
        c(m, g)
        a = m in ([s, u], not s and [t] * 3) and i(t, u) == a and i(t, s) or a
    h = m * 3 if m else [[s and [s], u, t], i(t, [s, u, t], 3) if s else u] if u else (a := i(*a, t, 2))
    k -= 1
    c(i(a, [k]), g + [i(h, (m or len(a) % 2 == m) and k)])
q = json.load(open(0)).pop
c(q(), [])
print('ACCEPT')
\end{lstlisting}

\subsection{Input, result and execution boundary}

Input is one JSON array. Its final element is the requested closed
formula; earlier four-field packets are consumed backwards. An unused
prefix is inert. A successful process exits zero and prints exactly
\texttt{ACCEPT}. An input or checking exception rejects with nonzero
exit status; the traceback is diagnostic. Merely finding the word
\texttt{ACCEPT} in arbitrary output is not the verdict convention.
There are no trusted Python \texttt{assert} statements, and the same
inference checks run under \texttt{-O}.

\subsection{Sorted array interpretation}

Canonical native formulas use
\[
       \D_{x,y}B\longmapsto[B],\qquad
       s=t\longmapsto[s,t],\qquad
       \Q_yB\longmapsto[B,B,B].
\]
Only the last field of a general ternary formula is its matrix; the
first two are annotations. Numerical terms use scalar zero, unary
successor and ternary scalar iteration. Both interpretation and
canonical encoding are sorted: the same array may have different
meanings at a numerical and a formula position.

\begin{center}
\begin{tabular}{cll}
\toprule
Arity & Numerical meaning & Formula meaning\\
\midrule
0 & $0$ & $0=0$\\
1 & $S$ of its child & $\D$ of its matrix\\
2 & $0$; annotations ignored & Equality of the two terms\\
3 & $I(n,a;v.h)$ & $\Q$ of the last field\\
$\geq4$ & $0$; children ignored & $0=0$; children ignored\\
\bottomrule
\end{tabular}
\end{center}

The scalar numerical coordinates $0,1,2$ all denote zero at the
numerical root depth three. Under $p$ protected numerical slots,
coordinates $3+2j,4+2j$ name slot $j$, counted from the innermost slot,
for $0\leq j<p$. Negative admitted coordinates name rigid numerical
parameters. Formula root depth is $2p$, and numerical root depth is
$3+2p$. Successfully admitted scalar padding at formula positions denotes
the fixed truth $0=0$.

In a unary numerical successor, its child is decoded under two dummy
zero slots. In a ternary iterator, counter and initializer are decoded
under one dummy zero slot, while the update is decoded under its genuine
state slot. In a paired formula the two slots are $y,x$; in a closure
the one slot is $y$. Thus the single child-depth update
\[
                    d\longmapsto d+5-\operatorname{arity}
\]
gives the required increments $4,3,2$ at arities $1,2,3$.
Canonical encoding shifts variables past the indicated dummy slots.
Rigid negative names do not shift. The named producer implements these
conventions independently of the kernel.

At larger arities the entire expression already denotes a constant.
Traversing an ignored annotation can reject it; a successful traversal
does not change that constant interpretation. In particular, zero or
negative depth increments inside such annotations do not authorize a
claim that all raw trees are idempotently revalidated. Only the semantic
opening statement at meaningful sorted roots is needed.

\subsection{Opening and freshness}

At global floor $k$, the opener preserves a scalar in the interval
$k\leq a<d$. At a numerical position an escaping scalar is the
distinguished substitution hole, and its image is checked at numerical
root depth three before insertion. A missing image at a used numerical
hole rejects. At a formula position scalar padding instead erases to
truth. Induction on a successfully traversed finite tree, using the
preceding table, proves capture-free semantic opening at every
meaningful position.

An inserted numerical image has no unbound positive coordinate at its
own root. Its free parameters are rigid negative names. Hence copying
it beneath extra protected slots does not capture its free parameters;
its internal binders retain their indices. Raw equality of successfully
checked expressions consequently implies equality of decoded meanings.
This statement does not identify different numerical programs merely
because they happen to have the same standard value.

Each non-hypothesis continuation decreases $k$. The integer new floor is
the designated fresh parameter. Previously checked relevant goals and
retained hypotheses mention no newly admitted name: their earlier
no-image checks would have rejected such an occurrence. If noninteger
JSON numbers occur, any additional names in the newly admitted interval
are separate fresh parameters, not aliases of its integer endpoint.
Only finitely many occur in a finite certificate. The known-zero-instance
argument substitutes zero for all of them when necessary.

\subsection{The three packet cases}

Write a packet as $[m,s,u,t]$. A truthy major $m$ takes priority. The
checker proves it in the old context before proceeding. If it is the
indicated equality $[s,u]$, and the current goal is the template opened
at $u$, the goal changes to its opening at $s$. For a canonical closure
$[t,t,t]$, the source field must denote the zero instance; the same
operation implements QE. A mismatch leaves the goal unchanged. Either
way a proved major is required, so a failed major does not become a
harmless no-op.

The exported local expression is $m*3$. A unary paired major exports
its three-copy closure, with the witness opened at the fresh name.
Every other nonempty array major exports a high-arity truth. Relevant
scalar majors cannot create arbitrary unproved consequences: their
only admitted formula interpretation is truth. The joint profile's
additional iterable aliases are covered by its inherited decoder.

For a false-valued $m$ and truthy $u$, the initializer selects
computation. The packet constructs
\[
 \begin{cases}
 [\,[s,u,t],u\,], &s\text{ false-valued},\\
 [\,[[s],u,t],\ i(t,[s,u,t],3)\,], &s\text{ truthy}.
 \end{cases}
\]
The first is a zero computation; the second constructs an actual
successor counter. When an update uses its state, filling it checks the
entire recursive source, including its update, at numerical root depth
three. This prevents an escaping update coordinate from being mistaken
for an extra state variable. If the update does not use its state,
the final generated equation still validates its meaningful contents.
The inherited zero-instance proof applies after setting newly admitted
names to zero.

When both $m$ and $u$ are false-valued, the packet is paired introduction.
The unary goal is opened at the supplied witness $t$, retaining the
index slot. Its successor instance becomes the next goal. Marker
$m=0$ retains the arbitrary-index instance only for a literal equation;
the case marker $m=[]$ retains the zero instance. Compound formulas get
the zero instance. The literal-equation restriction therefore remains
active. Exact unary unpacking rejects a binary equation as a paired
introduction; the undersupplied empty-goal case has the harmless truth
interpretation.

The last assignment shares construction of the local hypothesis across
these cases. Python evaluates only the selected branch of the conditional
expression. In the introduction branch the named assignment updates
the goal and hypothesis to the same opened matrix, just as the separate
assignments did before hoisting. The equality-transport Boolean
expression in the byte and preferred profiles likewise changes no
canonical proof: object and closure matrices are nonempty arrays. If
a noncanonical transported source is false-valued, the expression keeps
the original goal, imposing the stronger unchanged obligation.

\subsection{Why a generated hypothesis need not be universally true}

For fresh names $\vec b$ absent from $\Gamma,A$, the following rule is
admissible:
\[
 \vdash H(\vec0),\quad \Gamma,H(\vec b)\vdash A
 \quad\Longrightarrow\quad\Gamma\vdash A.
\]
Substitute zero for the fresh names in the second finite derivation,
renaming later eigenparameters first, and cut in the first proof.
No consistency or truth oracle is used.

This matters for the new dispatch, since arithmetic can now be selected
with the otherwise-progress marker $m=0$. For example the packet data
\[
                         m=0,\ s=0,\ u=3,\ t=0
\]
can generate, at the fresh parameter $b$, the decoded assumption
\[
                         H(b): I(0,0;v.0)=b.
\]
Its zero instance is a computation theorem; its instance at one is
false in the standard model. The retained goal is validated before
$b$ exists. The tests accept an independent valid continuation after
this supply and reject an attempted closed false consequence. The
soundness argument is the admissible fresh-name rule, not a mistaken
claim that every generated open hypothesis is uniformly true.

\subsection{Soundness, completeness and profiles}

\begin{theorem}[Checker correspondence]
Every successful finite run of the explicit-list checker has a decoded
conclusion derivable in the native iPRA presentation. Every finite
canonical native derivation has a finite certificate after the displayed
computation normalization, under the usual unbounded-resource
mathematical interpretation of the checker.
\end{theorem}
\begin{proof}
For soundness, induct on the finite successful checking trace.
Hypothesis closure uses semantic opening. The major branch uses the
inductively proved major, transport or QE when their raw match holds,
and DE or truth weakening for the export. Introduction uses the exact
unary shape, the fresh-name invariant and the correct case/progress
assumption. Computation uses its derivable zero instance and the fresh
discharge lemma. The final continuation opening realizes the displayed
successor or witness substitution. Hoisting the hypothesis construction
preserves the evaluated branch and its state changes.

For completeness, first replace each excluded computation leaf by the
three-node zero-certificate repair. Use canonical arity encodings,
fresh negative labels, and the inherited backward packet serialization
of Hyp, Tr, QE, DE and the two introductions. Every remaining computation
has a truthy initializer and, in its successor case, a truthy predecessor.
It therefore reaches the selected branch and generates its exact
computation equation. Hyp closes that continuation. All retained
contexts and conclusions are preserved by this finite translation.
\end{proof}

The byte profile's predicate \texttt{a*0 == []} distinguishes lists
among successfully tested JSON values. Strings and numbers produce
different zero products; mappings and null raise errors. It is the
inherited compact spelling of a list test, not proof-supplied execution.
The preferred explicit list predicate makes this intent clearer.

The joint profile tries iterable traversal and falls back to scalar
handling on an exception. It accepts some additional harmless iterable
aliases. Nonempty strings may recurse indefinitely in an ideal
unbounded-stack execution; actual CPython rejects through its recursion
limit and subsequent scalar type failure. Its simultaneous record must
therefore be read as a fixed-runtime comparison and canonical-proof
completeness statement, not an ideal all-JSON termination theorem.
The explicit-list and zero-product profiles avoid that malformed-string
qualification. All three remain subject to ordinary physical resource
limits on large genuine proofs, as do the supplied baseline programs.

% END theory/checker.tex


\section{Executable evidence and reproducibility}
\label{sec:evidence}


% BEGIN evidence/validation_table.tex
\begin{center}
\begin{tabular}{lr}
\toprule
Verification item & Result\\
\midrule
Named and adversarial certificate cases & 36\\
Complete CLI runs, three profiles and ordinary/\texttt{-O} & 216\\
Expected successful runs & 108\\
Expected rejected runs & 108\\
Verdict mismatches & 0\\
Frozen Python source byte strings remeasured & 370\\
Driver-only matched controls reconstructed & 7\\
\bottomrule
\end{tabular}
\end{center}

% END evidence/validation_table.tex


The positive certificates include the auxiliary zero itself, each
repaired arithmetic boundary, an update using its state, equality
transport through a closure, literal-zero existential witnesses, actual
equation induction using its induction hypothesis, and nested contexts.
The negative cases include compound induction, escaping names, false
equations, malformed updates, incorrect endpoints and an unproved major.
New tests distinguish the rejected old zero-initializer packet from a
successful proof of the same theorem through the zero certificate.

The producer is a separate untrusted program. It constructs named proof
trees, applies the simultaneous three-node normalization, and serializes
them without importing a checker. It is a producer for the supplied
examples and the displayed native proof-tree grammar, not a completed
automatic compiler from arbitrary conventional iPRA input. General
completeness is the written proof transformation.

Finite tests verify implementations and counterexamples; they do not
prove universal soundness or bi-interpretability. The written arguments
and the scoped independent review are included. There is no proof-assistant
certification in this package.

Run the complete offline reproduction with
\begin{lstlisting}
python3 reproduce.py
python3 build.py
\end{lstlisting}
The first command verifies the manifest, remeasures all frozen source
bytes and controls, regenerates the named certificates, and replays the
selected implementations in ordinary and optimized Python. The second
compiles the standalone LaTeX source and requires a normal TeX Live
installation. The archived proof-checker programs import only Python's
standard JSON module. Historical corpus sources are parsed for
measurement and are not executed by the corpus remeasurement.

\section{Two reductions that were rejected}
\label{sec:failed}

A smaller apparent branch count is not sufficient. We retain two
separate obstructions because they identify actual mathematical and
proof-completeness failures.

\subsection{Unconditional coercions lose reflexivity}

An experimental opener replaced every exposed numerical coordinate by
a binary projection box $[t,0]$. Numerically that box was intended as
identity; at a formula position it was interpreted as the zero test
$t=0$. This seemed to remove a sort-dependent branch. But the computation
rules supplied no projection equations.

For the precise archived experiment define a syntactic valuation $V$:
every nonbinary formula has value true, while
\[
 V([L,R])\quad\Longleftrightarrow\quad
 R=0\text{ literally}\ \text{or}\ L\text{ is a ternary array}.
\]
Every successful opening preserves $V$: zero stays zero, an inserted
box is binary rather than ternary, and existing array arities are
preserved. Computation supplies have a ternary iterator on the left;
major exports have arity divisible by three; equality transport and
closure specialization preserve $V$; paired introduction has a unary
conclusion. Induction on finite successful traces shows that every
accepted closed goal has value true.

The experiment represented $S0$ by $[0,0,0,0]$. Its ordinary reflexivity
equation $[S0,S0]$ has value false, so no finite certificate can establish
it. Adding an outer projection box to both endpoints has the same
problem. The experiment is not a delivered iPRA checker. This is a
completeness obstruction, not a claim of an arithmetically false accepted
theorem. It shows why a plausible semantic alias does not automatically
provide a usable internal equality.


% BEGIN theory/eager_obstruction.tex
% Original research note for the present investigation.  This fragment is
% self-contained modulo amsmath, amssymb, and amsthm.  It does not claim
% proof-assistant verification or literature-wide priority.

\section{A finite-model obstruction to eager initialization}
\label{sec:eager-obstruction}

There is a tempting way to shorten the computation clause of an iteration
checker.  Instead of preserving an independent initializer at counter zero,
apply the update immediately:
\begin{align}
 F(0,a;v.h)&=h[a/v],\label{eq:eager-zero}\\
 F(Sn,a;v.h)&=h[F(n,a;v.h)/v].\label{eq:eager-step}
\end{align}
The same syntactic opening operation then appears in both computation
clauses.  This subsection proves that, with the stated remaining
primitives, this change cannot present iPRA.  The obstruction applies to
relational interpretations and to finite-dimensional interpretations with
quotients, as well as to translations by terms.

\paragraph{The precise proposed theory.}
Let the numerical terms be the finite expressions
\[
 t ::= x\mid 0\mid S(t)\mid F(t,a;v.h).
\]
In the last constructor, $v$ is bound in the update $h$ only.  Every other
free variable of $h$ is a passive parameter.  Nested occurrences of $F$
are allowed, but this is a finite structural grammar: the update does not
refer circularly to the operation currently being defined.  Equivalently,
one may use a staged first-order signature containing a separate operation
symbol for each already constructed update description.  Substitution is
capture avoiding.

Let $T_{\mathrm e}$ use ordinary intuitionistic first-order logic with
actual equality, the equations \eqref{eq:eager-zero}--\eqref{eq:eager-step},
the two boundary assertions
\[
 \forall x\,(x=0\vee\exists y\,x=Sy),\qquad 0\ne S0,
\]
and induction on literal equations, with arbitrary retained assumptions
and passive parameters.  The countermodel below also satisfies the
stronger assertion $\forall x\,(0\ne Sx)$ and induction on \emph{every}
formula.  It therefore covers any version obtained by strengthening the
induction schema in that way.  Successor injectivity is \emph{not} an
axiom of this proposed simplification.

\paragraph{Two-element interpretation.}
Take the ordered set $M=\{0,1\}$ with $0<1$, interpret zero by $0$, and put
\[
                         S^M(b)=1\quad(b\in M).
\]
For a variable assignment $\rho$, interpret the eager constructor by
\[
 \llbracket F(n,a;v.h)\rrbracket_\rho
   :=\llbracket h\rrbracket_{\rho[v\mapsto\llbracket a\rrbracket_\rho]}.
 \tag{Eager interpretation}
\]
The counter is deliberately ignored.  This equation is a definition of
the mathematical interpretation, not a new evaluator trusted by the
proposed proof checker.

\paragraph{Monotonicity lemma.}
Every numerical term is monotone in all its free variables under this
interpretation.

\paragraph{Proof.}
Use structural induction on the finite term.  Variables are projections,
zero is constant, and successor is the constant function $1$.  For the
remaining case let
\[
 A(\rho)=\llbracket a\rrbracket_\rho,
 \qquad H(b,\rho)=\llbracket h\rrbracket_{\rho[v\mapsto b]}.
\]
After renaming the bound variable if necessary, the induction hypothesis
says that $A$ is monotone and $H$ is monotone jointly in $b$ and its passive
parameters.  If $\rho\le\sigma$ coordinatewise, then
\[
 H(A(\rho),\rho)\le H(A(\sigma),\sigma).
\]
This is exactly the assertion for $F(n,a;v.h)$.  Dependence on the
counter term may disappear, which causes no problem for monotonicity.
Nested eager binders are handled by the same structural induction.
\hfill$\square$

For completeness, the usual substitution lemma holds for this
interpretation: evaluating a capture-free term substitution is the same
as modifying the corresponding free-variable assignment.  Its proof is
another structural induction.  In the binder case first rename $v$ away
from the substituted variable and replacement term, and then apply the
induction hypotheses to $a$ and $h$.  Consequently the displayed
interpretation respects the formal substitution in both computation
schemas, including arbitrary fixed passive parameters.

\paragraph{Idempotence lemma.}
Every monotone map $H:M\to M$ is idempotent.

\paragraph{Proof.}
The possible pairs $(H(0),H(1))$ are precisely
\[
                       (0,0),\qquad(0,1),\qquad(1,1).
\]
They are respectively the constant-zero map, the identity, and the
constant-one map.  Each satisfies $H(H(b))=H(b)$ for both $b\in M$.
\hfill$\square$

\paragraph{Finite-model theorem.}
The structure $M$ is a model of $T_{\mathrm e}$, and remains a model if
literal-equation induction is replaced by full first-order induction.

\paragraph{Proof.}
Fix arbitrary values of all passive parameters and let $H$ be the
resulting unary interpretation of the update $h$.  By the monotonicity
lemma $H$ is monotone, and therefore idempotent.  If $a^M$ is the value of
the initializer, the zero clause reads
\[
                         H(a^M)=H(a^M).
\]
The left side of the successor clause is again $H(a^M)$, because the
counter is ignored; its right side is $H(H(a^M))$.  These agree by
idempotence.  The substitution lemma justifies these calculations for
the schemas as written, not merely for closed numerical instances.

Coverage holds because the elements are $0$ and $S0=1$.  Also $Sx=1\ne0$
for each $x$.  To check induction, fix any formula $A(x,\vec p)$ and an
assignment of its passive parameters.  If $A(0,\vec p)$ and
\[
             \forall x\,(A(x,\vec p)\to A(Sx,\vec p))
\]
hold, the step at $x=0$ gives $A(1,\vec p)$.  These are all elements of
the domain, so $\forall x\,A(x,\vec p)$ holds.  This argument works for
arbitrary formulas and under any fixed retained context.  The classical
structure is also a one-world intuitionistic Kripke model, so all
intuitionistic first-order rules with equality are valid in it.
\hfill$\square$

\paragraph{The native paired-binder presentation is covered.}
If the eager arithmetic is instead combined with the inherited
Eq/$Q$/$D$ presentation, interpret its clauses by
\[
 Q_yB(y):=B(0)\to\forall y\,B(y),\qquad
 D_{x,y}B(x,y):=\exists x\bigl(B(x,0)\to\forall y\,B(x,y)\bigr).
\]
Hypothesis use, equality transport, $Q$ specialization from its zero
instance, and fresh-witness $D$ elimination are sound in this ordinary
reading.  For case introduction, choose the supplied witness $t$.
The retained base assumption supplies the $y=0$ instance, while the
fresh-successor premise supplies the $y=1=S0$ instance.  For progress
introduction, the retained base assumption and the forward step imply
every instance by the induction just verified.  The usual freshness
conditions keep the witness and passive parameters fixed in these
arguments.  Thus neither removing ordinary connective primitives nor
packaging induction as the native progress rule removes the countermodel.

\paragraph{Corollary: no interpretation of iPRA.}
There is no ordinary finite-dimensional first-order interpretation of
iPRA in $T_{\mathrm e}$, even if its interpreted domain is a definable
subset of $M^d$ modulo a definable equivalence relation.  Consequently
$T_{\mathrm e}$ cannot satisfy the required bi-interpretability condition.

\paragraph{Proof.}
An interpretation would turn every model of $T_{\mathrm e}$ into a model
of iPRA.  In the displayed two-element model, a definable subset of $M^d$
has at most $2^d$ elements, and a quotient cannot increase that number.
A finite collection of interpreted sorts does not change finiteness.
But iPRA proves that all standard successor numerals are distinct, using
successor injectivity and zero--successor separation; hence it has no
finite model.  This is a contradiction.  The same argument permits a
fixed finite tuple of parameters in the interpretation.
\hfill$\square$

\paragraph{Scope and relation to the earlier obstructions.}
This result concerns eager initialization $h(a)$ with an independent
initializer $a$.  It is not the existing almost-unary growth argument
for a tied seed or a zero-seeded restricted iterator.  The supplied
\texttt{local\_scopes} draft also already excludes relational
bi-interpretability for an existential-free tied-seed theory, by finite
witness extraction and finite-profile image counting.  The present
argument obstructs all ordinary finite-dimensional relational
interpretations, including quotients, because it exhibits a finite
model of this different eager theory.  In particular it excludes even
a one-way interpretation of iPRA; it is not a claim to be the first
obstruction to relational bi-interpretability in this research program.
It does not apply to a recursion
update with a separately available counter argument, to the initialized
iterator $I(0,a;h)=a$, or to the successor-seeded iterator
$R(0,a;h)=Sa$.  Those alternatives require their own analyses.
Adding successor injectivity would exclude this countermodel, but would
be an additional arithmetic axiom and would not by itself prove a new
primitive-recursion compiler or bi-interpretation.  This obstruction
rules out the particular apparent simplification; it is not a lower
bound on all iPRA axiomatizations and is not a positive size improvement.

% END theory/eager_obstruction.tex


\section{Assessment and remaining limits}

The positive result is an exact admissibility theorem with an executable
consequence: a computation tag disappears, the complete source becomes
smaller, and all existing interpretation obligations are preserved by
the identity translation. The independent zero certificate is the
reason the syntactic restriction is legitimate. The shared assignment
and inherited representation profiles contribute further implementation
savings, which are separated from that theorem in the comparison.

The result does not establish an absolute simplest theory. It keeps the
inherited arithmetic language, iterator, and logical reconstruction.
Calling this a significant new arithmetic foundation would overstate
the finding. The strongest simultaneous source comparison also uses the
explicitly qualified generic traversal. A completely explicit-list
implementation below every matched token/AST/byte minimum remains a
stronger implementation target, even though its present byte size is
already substantially smaller.

The next formalization target is the decoder/opening invariant together
with known-zero-instance supply. The new normalization theorem is short
and independent of the difficult scalar predecessor bootstrap, making it
a convenient separately checkable contribution.

\clearpage
\appendix
\section{Attribution for the complete inherited bridge}

The following five appendices are reproduced from the arity-calculus
archive in the pinned snapshot, where they retain their earlier
provenance from generated-assumptions and scalar developments. Their
text is unchanged. Statements such as ``new'' within an inherited
appendix refer to its original report. They are supplied to make the
logical and arithmetic bridge available in this package; they are not
claimed as new results here. The exact input hashes and MIT No Attribution
license are included in \file{reference/}.


% BEGIN reference/inherited_choice_basis.tex
% New mathematical result: complementary introduction choices.
% Requires amsmath/amssymb/amsthm and theorem, lemma, proposition environments.
% The inherited arithmetic-to-iPRA bridge must be supplied separately.

\section{Complementary introductions with one local assumption}
\label{sec:choice-introduction}

The old paired introduction retains both the zero instance and an
induction instance when its matrix is an equation. Both together are
unnecessary. A certificate may choose either one, provided the induction
choice is restricted to a literal equation. This changes neither the
formula language nor the meaning of a formula. The omitted combination
is recovered by a short nested introduction--elimination construction.

\subsection{Complete native syntax and rule inventory}

The scalar numerical language is
\[
 t::=x\mid0\mid St\mid I(n,a;u.h).
\]
Only $h$ binds the state variable $u$; $n$ and $a$ do not.
The update may contain fixed numerical parameters. The computation
axioms in every retained context are
\[
 I(0,a;u.h)=a,\qquad
 I(Sn,a;u.h)=h[I(n,a;u.h)/u].
 \tag{I0, IS}\label{eq:choice-computations}
\]
The argument also applies unchanged to the full fixed-description PR
language with its projection, composition, and recursion equations.
No numerical variable ranges over function descriptions, and no uniform
evaluator is postulated.

Native formulas and reusable context clauses are
\[
 A::=(s=t)\mid\mathsf D_{x,y}A,\qquad
 \langle A;y.B(y)\rangle,\quad y\notin\operatorname{FV}(A).
\]
Bound variables are renamed before capture-free substitution.
The two variables $x,y$ in a paired formula are distinct.
An ordinary assumption $F$ is the constant clause
$\langle H;v.F\rangle$, where $v$ is absent from $F$ and $H$ is the
fixed computation theorem $I(0,0;u.0)=0$. Thus ordinary hypothesis
use is a clause Use preceded by this computation proof.

There are computation axioms and exactly the following inference
patterns:
\[
\frac{\Gamma\vdash A}{\Gamma\vdash B(t)}
\quad\mathsf{Use}
\quad(\langle A;y.B(y)\rangle\in\Gamma),
\]
\[
\frac{\Gamma\vdash\mathsf D_{x,y}B(x,y)\qquad
      \Gamma,\langle B(a,0);y.B(a,y)\rangle\vdash F}
     {\Gamma\vdash F}
\quad\mathsf{DE},
\qquad
\frac{\Gamma\vdash s=t\qquad\Gamma\vdash F(s)}
     {\Gamma\vdash F(t)}
\quad\mathsf{Tr},
\]
\[
\frac{\Gamma,B(t,\eta_\varepsilon(z))\vdash B(t,Sz)}
     {\Gamma\vdash\mathsf D_{x,y}B(x,y)}
\quad\mathsf{DI}_{\varepsilon},
\qquad
\eta_0(z)=0,\quad\eta_1(z)=z.
\tag{Choice DI}\label{eq:choice-di}
\]
Here $\varepsilon$ is a certificate choice in $\{0,1\}$, and
$\varepsilon=1$ is permitted only when $B$ has a literal equation
root. An equation beneath a binder is not a literal equation.
Call $\varepsilon=0$ the \emph{case choice} and $\varepsilon=1$
the \emph{progress choice}. Both retain exactly one ordinary
formula assumption. Compound matrices have only the case choice.

In DI the eigenparameter $z$ is fresh for $\Gamma$, the witness
$t$, and every fixed parameter of $B$. In DE the eigenparameter
$a$ is fresh for $\Gamma$, the conclusion $F$, and every fixed
parameter of $B$. In transport, $F(v)$ is any fixed native
formula template. Displayed local assumptions and clauses are
discharged. There are no primitive reflexivity, coverage,
separation, ordinary logical, or additional induction rules.

We call this calculus $\mathcal C$. The original comparison
calculus $\mathcal N$ has identical syntax, computation,
Use, DE, and transport, but its introduction is
\[
\frac{\Gamma,B(t,0),[B(t,z)]_{\mathrm{at}}\vdash B(t,Sz)}
     {\Gamma\vdash\mathsf D_{x,y}B(x,y)}.
\tag{Old DI}\label{eq:choice-old-di}
\]
The bracket supplies $B(t,z)$ exactly at a literal equation.

\subsection{Exact equivalence and the unused-witness construction}

\begin{lemma}[Structural admissibility]
\label{lem:choice-structural}
Weakening, exchange, contraction, renaming, ordinary-assumption cut,
and capture-free numerical substitution are admissible in $\mathcal C$.
\end{lemma}
\begin{proof}
Induct on derivations. Rename local eigenparameters before inserting
proofs or terms. Numerical substitution preserves a formula's root
constructor, so an allowed progress choice remains allowed.
The selected choice does not change under substitution or cut.
For ordinary-assumption cut, replace each use of its constant clause
by the supplied proof. Computation schemas are closed under
capture-free numerical substitution. These facts verify every case.
\end{proof}

\begin{theorem}[Identity-language equivalence]
\label{thm:choice-equivalence}
For every retained clause context $\Gamma$ and native formula $F$,
\[
 \Gamma\vdash_{\mathcal C}F
 \quad\Longleftrightarrow\quad
 \Gamma\vdash_{\mathcal N}F.
\]
The transformations are effective on finite derivations and leave
all numerical and formula syntax unchanged.
\end{theorem}
\begin{proof}
Each new introduction is an instance of the old one after weakening
by the omitted local assumption. For the case choice at an equation,
add the induction instance. For the progress choice, add the zero
instance. Compound case introduction is already the old rule.
Thus every new derivation translates to an old one.

For the reverse direction, transform an old derivation recursively.
Only its atomic DI case needs a construction. Suppose its matrix is
the equation $E(x,y)$ and its transformed premise is
\[
 \pi:\quad\Gamma,E(t,0),E(t,z)\vdash E(t,Sz).
\]
The outer new inference will be case DI with witness $t$. It is
therefore enough to prove $E(t,Sz)$ under
$\Gamma,E(t,0)$.

Rename the eigenvariable in $\pi$ to a fresh $v$. Choose a new
variable $u$ absent from $E(t,y)$ and all retained data. Under
$\Gamma,E(t,0)$, progress DI with witness zero introduces
\[
             \mathsf D_{u,y}E(t,y)
\]
from exactly the renamed premise $\pi[v/z]$.
The bound coordinate $u$ is deliberately unused.

Apply DE to this paired formula. Because $u$ is unused, the opened
clause is exactly
\[
               \langle E(t,0);y.E(t,y)\rangle;
\]
it contains no new existential witness. The retained ordinary
assumption $E(t,0)$ supplies its antecedent. Use at $Sz$ gives
$E(t,Sz)$, and DE closes. Finally the outer case DI introduces
the original $\mathsf D_{x,y}E(x,y)$ with witness $t$.
The original premise proof $\pi$ has been inserted once.

An old compound introduction already is the new case choice.
All other rule cases are unchanged. This completes the induction
on the original derivation.
\end{proof}

The unused inner coordinate is essential. Introducing
$\mathsf D_{x,y}E(x,y)$ internally and then eliminating it would
give an arbitrary fresh witness, not necessarily the original
$t$. The construction instead keeps $t$ as a fixed parameter
of $E(t,y)$ and binds only the unused coordinate $u$.

\begin{proposition}[Linear bound on inference nodes]
\label{prop:choice-node-bound}
Count each computation axiom, Use, DI, DE, and transport as one
inference node, expanding ordinary-assumption use as a computation
of $H$ followed by Use. If an old derivation has $N$ nodes and
$r$ atomic DI nodes, the displayed transformation has exactly
\[
                   N+5r\leq 6N
\]
nodes.
\end{proposition}
\begin{proof}
Each atomic old DI node is replaced by six nodes: outer case DI,
inner progress DI, DE, Use of the opened clause, Use of the
retained base assumption, and the computation of $H$ needed for
that last ordinary-assumption use. Its old premise subtree
occurs once. All other inference nodes are copied once.
\end{proof}

This counts inference-tree nodes. It is not a bound on serialized
JSON bytes, formula sizes, checker runtime, or the number of
consumed packet cells when hypothesis closure is implicit.
In the formula-context Q/D presentation, the corresponding
replacement has outer case DI, inner progress DI, DE, QE,
and two Hyp nodes, giving the same six-node count.

\subsection{Induction and numerical cases separate cleanly}

\begin{lemma}[Contextual equation induction]
\label{lem:choice-ei}
The progress choice, DE, and Use derive
\[
\frac{\Gamma\vdash E(0)\qquad
      \Gamma,E(z)\vdash E(Sz)}
     {\Gamma\vdash E(t)}
\tag{EI}\label{eq:choice-ei}
\]
for every equation $E$, with the usual freshness condition.
\end{lemma}
\begin{proof}
Choose an unused witness coordinate $u$. Progress DI introduces
$\mathsf D_{u,y}E(y)$ from the given step. DE opens the clause
$\langle E(0);y.E(y)\rangle$. Supply its antecedent using the given
base proof and use the clause at $t$. Close DE.
\end{proof}

\begin{lemma}[Cases, equality, and binder congruence]
\label{lem:choice-preinduction}
The case choice derives, for arbitrary native formulas,
\[
\frac{\Gamma\vdash F(0)\qquad\Gamma\vdash F(Sz)}
     {\Gamma\vdash F(t)}
\tag{NC}\label{eq:choice-nc}
\]
with the usual freshness condition. Equality reflexivity,
symmetry, transitivity, numerical congruence, and uniform
matrix interderivability under $\mathsf D$ are derivable
without using the progress choice.
\end{lemma}
\begin{proof}
For cases, use case DI with an unused witness coordinate and
matrix $F(y)$, ignoring its retained base in the successor
premise. DE opens $\langle F(0);y.F(y)\rangle$; the given zero
proof and Use at $t$ finish.

For reflexivity put $r=I(0,t;u.0)$. Use the computation equation
$r=t$ as both transport premises with motive $v=t$ to obtain
$t=t$. The usual explicit equality motives give symmetry,
transitivity, and numerical congruence.

For binder congruence, open the old paired formula at a fresh
witness $a$ and introduce the new one by case DI at $a$.
Convert its retained base to the old base, use the opened
old clause at the successor, and convert back. Ordinary cut
closes the construction, and the reverse conversion gives
the converse. No induction is used.
\end{proof}

The primitive progress choice is still restricted to equations.
Allowing it at arbitrary compound matrices would, by the
argument for \eqref{eq:choice-ei}, supply induction for
arbitrary formulas and change the intended theory.

\subsection{The entire ordinary logical interface}
\label{subsec:choice-logic}

Here is a predecessor-free reconstruction of the ordinary logical
operations. It makes explicit that the case choice supplies the
logical interface independently of progress induction.

Let $C(v,a,b):=I(v,a;u.b)$ with $u$ unused. Computation gives
$C(0,a,b)=a$ and $C(Sv,a,b)=b$. Define the gate on fixed finite
formula syntax by
\[
 G_v(s=t):=(C(v,0,s)=C(v,0,t)),\qquad
 G_v(\mathsf D_{x,y}B):=\mathsf D_{x,y}G_vB,
\]
renaming bound variables away from the control. Equality,
binder congruence, and external structural induction give
\[
 \vdash G_0A,\qquad G_{Sv}A\dashv\vdash A,\qquad
 A\dashv\vdash\mathsf D_{u,y}G_yA
 \quad(u,y\notin\operatorname{FV}(A)).
\tag{Gate and padding}\label{eq:choice-gate}
\]
For padding, case DI proves the forward direction. DE, the
zero-gate theorem, and Use at $1$ prove the reverse direction.

Define a selector between two arbitrary formulas by
\[
 M_v(s=t,u=w):=(C(v,s,u)=C(v,t,w))
\]
and, with aligned fresh paired variables, by
\[
\begin{aligned}
 M_v(\mathsf D_{x,y}A,\mathsf D_{x,y}B)
   &:=\mathsf D_{x,y}M_v(A,B),\\
 M_v(E,\mathsf D_{x,y}B)
   &:=\mathsf D_{x,y}M_v(G_yE,B),\\
 M_v(\mathsf D_{x,y}A,E)
   &:=\mathsf D_{x,y}M_v(A,G_yE),
\end{aligned}
\]
where $E$ is an equation. Every call decreases maximum input
binder depth. Padding and congruence prove
$M_0(A,B)\dashv\vdash A$ and
$M_{Sv}(A,B)\dashv\vdash B$.

Put $\delta(v)=C(v,1,0)$ and $p(z)=I(z,0;u.\delta(u))$.
Only the fixed computations $p(1)=1$ and $p(2)=0$ are needed.
The full logical macro inventory is
\[
\begin{aligned}
 A\Rightarrow_*B&:=\mathsf D_{u,z}M_z(A,B),\\
 A\wedge_*B&:=\mathsf D_{u,z}G_zM_{p(z)}(A,B),\\
 \exists_*x\,A(x)&:=\mathsf D_{x,y}G_yA(x),\\
 A\vee_*B&:=\exists_*x\,M_x(A,B),\\
 \forall_*x\,A(x)&:=
        A(0)\wedge_*\mathsf D_{u,z}G_zA(z),\\
 \bot_*&:=(0=1).
\end{aligned}
\tag{Logical macros}\label{eq:choice-macros}
\]
Newly bound variables are fresh; $u$ is unused.

Implication introduction obtains its assumption from the zero
selector and supplies the consequent at successors, then uses
case DI. Elimination opens its clause and uses it at $1$.
Existential introduction uses its witness and the successor
gate; elimination opens at a fresh witness, supplies the
zero-gate theorem, and uses the clause at $1$.

Disjunction introduction uses witnesses $0$ and $1$.
For elimination, discharge the two branch proofs as
$A\Rightarrow_*F$ and $B\Rightarrow_*F$. NC applied to
$M_z(A,B)\Rightarrow_*F$ gives this implication at every
control: the two endpoint cases follow from the branch
implications. Existential elimination provides $M_a(A,B)$
and the implication gives $F$.

Given $A$ and $B$, selector endpoints and NC give $M_v(A,B)$
at an arbitrary control. Substitute $v=p(Sz)$ and use the
successor gate for conjunction introduction. For projections,
open its zero-gated clause and specialize at $1$ and $2$;
the two fixed computations give $B$ and $A$, respectively.

These rules prove existential numerical coverage:
\[
     (x=0)\vee_*\exists_*v\,(x=Sv).
\tag{Coverage}\label{eq:choice-coverage}
\]
Use NC. Zero takes the left injection of reflexivity, and a
successor takes the right injection with the evident witness.
Coverage gives equality cases even when the retained context
contains $x$; its elimination adds $x=0$ or $x=Sv$ rather
than silently changing that retained context.

Universal introduction follows from a fresh-parameter proof
by substituting zero and successors and forming the two
conjuncts. For elimination, take its $A(0)$ conjunct and
open the paired conjunct. Its base is a zero-gate theorem.
Coverage of the target chooses the zero instance from
$A(0)$ or a successor instance from the opened clause
and gate removal.

Finally, under $0=1$, computation and transport give
$s=C(0,s,t)=C(1,s,t)=t$ for arbitrary numerical terms.
Every equation follows, and external induction on formula
syntax with case DI gives every native formula. This
supplies falsity elimination. No excluded middle for
quantified formulas has entered the construction.

\subsection{Whole-domain realizations and both formula composites}

For a standard first-order reading, let $T_{\rm eq}$ have
intuitionistic equality logic, the same numerical signature
and computation equations, numerical coverage, separation
$(0=1)\to\bot$, and contextual equation induction.
Its interpretation of the native syntax is
\[
 (s=t)^\circ=(s=t),\qquad
 (\mathsf D_{x,y}B)^\circ
   =\exists x\bigl(B^\circ(x,0)\to\forall y B^\circ(x,y)\bigr),
\]
\[
 \langle A;y.B\rangle^\circ
       =A^\circ\to\forall y B^\circ.
\]
The domain is the entire numerical domain and interpreted
equality is actual equality.

Case DI is sound in this realization by numerical coverage.
Progress DI is sound by contextual equation induction under
the conditional base. Use, DE, transport, and computation
have their displayed ordinary meanings. Conversely, the
native macros derive every ordinary logical rule,
\eqref{eq:choice-ei} derives equation induction, and
\eqref{eq:choice-coverage} derives coverage. Separation
translates to the identity implication
$(0=1)\Rightarrow_*(0=1)$.

For the essential paired formula composite,
\[
 \mathsf D_{x,y}B
 \dashv\vdash
 \exists_*x\bigl(B(x,0)\Rightarrow_*\forall_*y\,B(x,y)\bigr).
\tag{Paired composite}\label{eq:choice-paired-composite}
\]
Forwards, DE opens its clause at a fresh $a$; under $B(a,0)$
that clause supplies $B(a,y)$ at a fresh arbitrary $y$.
Introduce the universal, implication, and existential
macros and close DE. Backwards, eliminate the existential
at a fresh $a$. Introduce the paired formula with witness
$a$ by case DI: the retained base, supplied implication,
and universal elimination at $Sz$ give its successor premise.

The ordinary interpretations of the six logical macros
have their ordinary round trips. For conjunction, its true
base and the instances $1,2$ recover both operands;
the converse uses selector cases. The universal has all
successor instances in its paired conjunct and zero in
its first conjunct; coverage gives exactly the ordinary
universal. Implication, disjunction, and existential
quantification use their selector or gate endpoints.
Falsity uses the explicit ordinary separation bridge.
External induction on finite formulas, together with
\eqref{eq:choice-paired-composite}, gives both formula
composites. Clauses correspond in both directions to
$A\Rightarrow_*\forall_*yB$ by the derived logical rules.

The numerical composites are literal identities with
internally provable graph $x=y$. No standard-numeral
restriction, retract, quotient, or equality reflection
assumption is being inserted at this stage. The separately
supplied arithmetic bridge identifies $T_{\rm eq}$ with
full iPRA; the present change preserves that bridge
unchanged because Theorem~\ref{thm:choice-equivalence}
preserves every native sequent.

\subsection{What the theorem does and does not claim}

The new inference interface retains one ordinary assumption
at every paired introduction. It preserves the entire
consequence relation, not just equations or closed
theorems. Consequently the established full-iPRA
interpretations and equiconsistency transfer unchanged.

This result is different from retaining only progress
introduction and strengthening closure elimination.
Here the original elimination rule is retained, and
the case and progress introductions coexist. It is also
different from allowing arbitrary compound progress,
which would supply stronger induction.

A complete-checker comparison must still count the
certificate choice, its equation-root restriction,
substitution, context handling, computation, and the
public input/output interface. The mathematical theorem
alone is not a byte-count or minimality theorem.
The inference-node bound is distinct from source-size
measurements and from serialized certificate size.
The proof is a written universal proof; finite regression
tests do not constitute a Lean/Rocq formalization of it.

% END reference/inherited_choice_basis.tex


% BEGIN reference/inherited_ordinary_bridge.tex
\section{Arithmetic strength and both ordinary interpretation composites}
\label{sec:ordinary-bridge}

We spell out the arithmetic bridge to avoid identifying standard numerical
agreement with bi-interpretability. First elaborate every fixed iterator
into an ordinary function symbol. At stage zero the symbols are $0,S$.
Given a fixed earlier-stage term $h(u,\vec p)$, introduce
$i_h(n,a,\vec p)$ with the two equations
\[
i_h(0,a,\vec p)=a,\qquad
i_h(Sn,a,\vec p)=h(i_h(n,a,\vec p),\vec p).
\]
Only finite descriptions occur; the domain contains no code for a
universal evaluator. Formal parameters are distinct when a description
is declared, while actual arguments can coincide.

Let $T_I$ be ordinary intuitionistic first-order logic with actual
equality and genuine falsity, these staged equations, contextual
induction for literal equations, and
\[
\forall x\,(x=0\vee\exists y\,x=Sy),
\qquad (0=S0)\to\bot.
\]
The second sentence is the falsity bridge. It must be stated: positive
numerical equations by themselves admit a singleton structure. Native
explosion from $0=S0$ and a translation of ordinary falsity to that
equation do not authorize silently dropping the bridge from $T_I$.

The native-to-ordinary translation is
\[
(s=t)^\circ=(s=t),\qquad
(\mathsf D_{x,y}B)^\circ=
\exists x\bigl(B^\circ(x,0)\to\forall y B^\circ(x,y)\bigr).
\]
A clause $\langle A;y.B\rangle$ becomes $A^\circ\to\forall yB^\circ$.
For the auxiliary closure used by the compact checkers,
$(\mathsf Q_yB)^\circ=B^\circ(0)\to\forall yB^\circ$.
The interpretation domain is all elements, defined by $x=x$;
equality is actual equality.

Case DI is sound because numerical coverage combines its retained base
with its successor proof. Progress DI is sound by literal-equation
induction under a temporarily assumed base. Use, DE and transport are
ordinary logical inferences under the translations. Conversely the
derived logical operations, equation induction, cases and explosion in
Section~\ref{sec:choice-introduction} recover the rules and axioms of
$T_I$. An auxiliary closure is interderivable with
$\mathsf D_{v,y}B(y)$ for unused $v$: in one direction use case DI with
the closure and its base; in the other eliminate the paired formula
and return its supplied closure. Thus formulas and context clauses have
the stated meanings in both presentations.

\subsection{The arithmetic dependency order}

The scalar update has no recursion-counter argument. We do not conceal
one in an encoding. The inherited internal proofs, reproduced in the
appendices, establish the required arithmetic in this order:
\begin{enumerate}
\item Using equation induction and numerical cases, construct a fixed
scalar predecessor term and prove $P0=0$ and $P(Sn)=n$.
Existential coverage then gives $x=0\vee x=S(Px)$: in the successor
case $x=Sy$, apply $P$ to obtain $Px=y$ and substitute.
\item Construct truncated subtraction and a numerical characteristic
of actual equality. Prove its reflection on the entire original domain.
\item Compile each fixed quantifier-free formula into its characteristic
equation and recover every contextual quantifier-free induction instance.
\item With that induction available, construct pairing and projections
by a bounded square-root search and compile every fixed general
primitive-recursive description into a scalar iterator term.
\end{enumerate}
No pairing or faithful equality detector is assumed to obtain
predecessor. No appeal to standard-model correctness replaces an
internal equality. Every induction fixes its passive parameters;
diagonal substitution occurs only into a previously proved identity.

For a target recursion
\[
f(0,\vec p)=g(\vec p),\qquad
f(Sn,\vec p)=h(n,f(n,\vec p),\vec p),
\]
let $\pi,L,H$ be the proved pairing and projections. The compiler uses
\[
V(n,\vec p)=I\bigl(n,\pi(0,g^*(\vec p));
u.\pi(SL(u),h^*(L(u),H(u),\vec p))\bigr),
\qquad f^*(n,\vec p)=H(V(n,\vec p)).
\]
Equation induction proves $L(V(n,\vec p))=n$; the two projection
identities then prove the target defining equations. The construction
is external recursion on each finite PR description, not an internal
quantification over arbitrary functions.

\begin{theorem}[Whole-domain bi-interpretability and equiconsistency]
The ordinary realization $T_I$ and full first-order iPRA are
bi-interpretable. In both theories $x=y$ defines the isomorphism from
the original domain to its double interpretation. The native Choice-DI
calculus is therefore a faithful proof presentation of a theory
bi-interpretable and equiconsistent with iPRA.
\end{theorem}
\begin{proof}
Interpret an iterator in iPRA by primitive recursion that ignores its
counter argument in the update. iPRA proves its equations, literal
equation induction, coverage and separation. In the reverse direction
use the scalar compiler. The predecessor and equality-reflection proofs
supply the successor axioms and every quantifier-free induction instance.
All ordinary logical rules are preserved.

For two iterators with the same seed and provably equal updates,
equation induction on the equality of their values proves uniqueness at
every counter, with all parameters fixed. The corresponding argument
applies to ordinary primitive recursion. Induct externally on a finite
description: its double translation satisfies its original defining
equations, so internal uniqueness identifies the two functions. This
also proves coherence when substitution changes the staged syntactic
description of an update.

The domain and equality are unchanged; hence the identity graph is
internally total, bijective and preserves every interpreted operation.
The displayed ordinary meanings of the logical macros and the paired
binder give both formula composites by external induction on finite
formula syntax. Clauses are similarly equivalent to their conditional
universal meanings. The explicit falsity bridge preserves contradiction
in both directions. Effective translations of finite proofs give
equiconsistency, without a claim that either theory proves its own
consistency.
\end{proof}

% END reference/inherited_ordinary_bridge.tex


% BEGIN reference/scalar_predecessor.tex
% Adapted with attribution from the pinned unparenthesized one-hole archive.
\section{Internal scalar predecessor bootstrap}
\label{sec:hybrid-bootstrap}

This scalar bootstrap is inherited from the supplied scalar development~\cite{erasure} and is reproduced to keep the scalar checker self-contained. It is not a new contribution of this paper. The proof uses a factor map fixing odd states and a saturated-iteration construction of predecessor.

The available arithmetic terms are $0$, $S$, and
\[
 I(n,a;u.h),\qquad
 I(0,a;u.h)=a,\qquad
 I(Sn,a;u.h)=h[I(n,a;u.h)/u].
\]
The update binds only $u$.  All other variables in it are fixed numerical
parameters.  We use equality reasoning, the already reconstructed
zero/successor cases, and induction on a single literal equation.
Every induction below keeps its other parameters fixed.  Predecessor,
successor injectivity, pairing, and induction on arbitrary
quantifier-free formulas are not assumptions of the argument.

\subsection{Elementary terms and two finite cycles}

Write $1=S0$ and use larger numerals as successor abbreviations.  Define
\[
 C(n,a,b)=I(n,a;u.b),\qquad
 Z(n)=C(n,1,0),\qquad
 a+b=I(b,a;u.Su).
\]
The update variable of $C$ is absent from its body.  Consequently
\[
 C(0,a,b)=a,\qquad C(Sn,a,b)=b.
\]
Numerical cases show that $Z(n)$ is $0$ or $1$.
The usual commutative-monoid equations for addition follow by literal
equation induction.  For example, $0+b=b$ and $Sa+b=S(a+b)$ are
proved by induction on $b$; these yield commutativity by a further
equation induction, and associativity follows by induction on the last
summand.  We use these equations to reorder finite sums.  The notation
$2a$ means $a+a$ and $3a$ means $a+a+a$; no multiplication operation
is needed for this bootstrap.

Put
\[
 N(n)=I(n,0;u.Z(u)),\qquad
 j(u)=Z(u)+2N(u),\qquad
 M(n)=I(n,0;u.j(u)).
\]
The range of $N$ is contained in $\{0,1\}$: its zero value is $0$,
and its successor value is a value of $Z$.  This uses numerical cases
on its counter, not induction on a disjunction.  Hence
\begin{equation}
 N(n)+Z(N(n))=1,\qquad
 Z(Z(N(n)))=N(n),\qquad N(n+2)=N(n).
 \label{eq:hybrid-parity}
\end{equation}
Here the last equation follows from two successor computations and
the preceding equation.  Literal-equation inductions give
\begin{equation}
 N(2n)=0,\qquad N(2n+1)=1,\qquad N(3n)=N(n).
 \label{eq:hybrid-parity-scaling}
\end{equation}
For the first equation the successor step adds two to the argument of
$N$; the second follows by one computation.  In the third, increasing
$n$ adds three to the left argument and therefore applies $Z$ once to
both sides of the induction equation.

Finite computations now give
\[
 j(0)=1,\qquad j(1)=2,\qquad j(2)=0.
\]
The proof below needs just these three values and
$M(Sn)=j(M(n))$; it does not assume a division algorithm or a theorem
characterizing $M$ as remainder modulo three.

\subsection{A ramp with a global scaling identity}

Define
\[
 G(n)=I(n,0;u.u+M(u)+1).
\]
The central invariant is
\begin{equation}
 M(G(n))=N(n).
 \label{eq:hybrid-ramp-invariant}
\end{equation}
Its zero case is immediate.  For its successor step, retain the
single induction equation and split the Boolean value $N(n)$.
If it is $0$, the next $G$ state is $G(n)+1$; its $M$ value is
$j(0)=1$, equal to $N(Sn)$.  If it is $1$, the next $G$ state
is $G(n)+2$; two computations of $M$ change its value from $1$
to $2$ and then to $0$, again equal to $N(Sn)$.  This proves
the invariant by literal-equation induction.

It follows that
\[
 G(Sn)=G(n)+N(n)+1,\qquad G(n+2)=G(n)+3.
\]
Equation induction therefore yields
\begin{equation}
 G(2n)=3n.
 \label{eq:hybrid-ramp-double}
\end{equation}
A separate literal-equation induction proves
\begin{equation}
 G(3n)=3G(n)+N(n).
 \label{eq:hybrid-ramp-triple}
\end{equation}
For completeness, its successor calculation is
\begin{align*}
 G(3Sn)
 &=G(3n)+N(3n)+4
   =3G(n)+2N(n)+4,\\
 3G(Sn)+N(Sn)
 &=3G(n)+3N(n)+3+Z(N(n)).
\end{align*}
The expressions agree in both Boolean cases for $N(n)$.
The base is $0=0$.  In particular,
\eqref{eq:hybrid-ramp-triple} was proved internally; it was not
inferred from an intended floor-function interpretation of $G$.

Set
\[
 H(u)=C(N(u),G(u),u).
\]
The crucial choice is to retain $u$ when its parity is $1$.
The preceding equations give
\begin{equation}
 H(2u)=3u,\qquad H(1)=1,\qquad H(3u)=3H(u).
 \label{eq:hybrid-factor-map}
\end{equation}
To check the last identity, use $N(3u)=N(u)$ and split the Boolean
$N(u)$.  At $0$, both sides are $3G(u)$ by
\eqref{eq:hybrid-ramp-triple}; at $1$, both are $3u$.
This argument does not require that an arbitrary number has already
been represented as $2v$ or $2v+1$.

\subsection{A comparator from scalar iteration}

For every fixed update $h$, equation induction on $n$, keeping $m,a$
and all parameters fixed, proves the semigroup identity
\[
 I(m+n,a;h)=I(n,I(m,a;h);h).
\]
The base is the zero computation and the step applies $h$ to the
induction equation. This identity will also be used in the later
pairing construction.


For a fixed update $h$, equation induction on $y$, with its initializer
$a$ and all parameters of $h$ fixed, proves the rotation identity
\begin{equation}
 I(Sy,a;h)=I(y,h(a);h).
 \label{eq:hybrid-rotation}
\end{equation}
Both sides have base $h(a)$, and their successor equations apply the
same $h$ to the previous equality.  Here and below $I(y,a;h)$ is
an abbreviation for one fixed finite scalar term, not quantification
over functions.

By \eqref{eq:hybrid-factor-map}, another equation induction gives
\begin{equation}
 I(y,3a;H)=3I(y,a;H).
 \label{eq:hybrid-iterator-scaling}
\end{equation}
Its base is $3a=3a$.  In the step the fixed-parameter induction
equation rewrites the new left side as
\[
 H(3I(y,a;H))=3H(I(y,a;H)).
\]
The initializer parameter stays $a$ throughout this induction.
Equation induction also gives $I(y,1;H)=1$, since $H(1)=1$.

Define
\[
 E(x)=I(x,1;u.2u),\qquad
 L(x,y)=N\bigl(I(y,E(x);H)\bigr).
\]
Then the following entire recursion specification is provable:
\begin{equation}
 L(0,y)=1,\qquad L(Sx,0)=0,\qquad L(Sx,Sy)=L(x,y).
 \label{eq:hybrid-comparator}
\end{equation}
The first equation uses $E(0)=1$ and the fixed-point identity for
$H$.  The second uses $E(Sx)=2E(x)$ and $N(2a)=0$.
For the third,
\begin{align*}
 L(Sx,Sy)
 &=N\bigl(I(Sy,2E(x);H)\bigr)\\
 &=N\bigl(I(y,H(2E(x));H)\bigr)\\
 &=N\bigl(I(y,3E(x);H)\bigr)\\
 &=N\bigl(3I(y,E(x);H)\bigr)\\
 &=N\bigl(I(y,E(x);H)\bigr)=L(x,y).
\end{align*}
The successive steps use rotation, $H(2a)=3a$, iterator scaling,
and $N(3a)=N(a)$.  Thus $L$ has the comparison equations at all
arguments internally.  Its standard-model meaning as the indicator
of $x\le y$ is motivation, not a premise of the following argument.

\subsection{Saturated iteration gives predecessor}

We now need only the three equations
\eqref{eq:hybrid-comparator}.  Define
\[
 B_m(u)=C(L(Su,m),u,Su),\qquad
 W(t,m)=I(t,0;u.B_m(u)).
\]
The subscript in $B_m$ records a fixed parameter; $B_m$ is an open
term abbreviation, not a function variable.

Numerical cases on $e$ prove the selector equation
\[
 C(e,Sa,Sb)=S C(e,a,b).
\]
Equivalently one may use equation induction on $e$, whose induction
hypothesis is unused.  Therefore
\begin{align}
 B_{Sm}(Su)
 &=C(L(SSu,Sm),Su,SSu)\notag\\
 &=C(L(Su,m),Su,SSu)=S B_m(u),
 \label{eq:hybrid-conjugacy}\\
 B_{Sm}(0)&=C(L(S0,Sm),0,S0)=C(L(0,m),0,S0)=S0.
 \label{eq:hybrid-conjugacy-base}
\end{align}

Literal-equation induction on $t$, with $m$ fixed, proves
\begin{equation}
 W(St,Sm)=S W(t,m).
 \label{eq:hybrid-saturated-shift}
\end{equation}
The base is \eqref{eq:hybrid-conjugacy-base}.  The successor step is
\begin{align*}
 W(SSt,Sm)
 &=B_{Sm}(W(St,Sm))\\
 &=B_{Sm}(S W(t,m))\\
 &=S B_m(W(t,m))=S W(St,m).
\end{align*}
Only the second line uses the current induction hypothesis, at the
same fixed $m$.

Two subsequent equation inductions give
\begin{equation}
 L(Sn,n)=0,\qquad W(n,n)=n.
 \label{eq:hybrid-diagonals}
\end{equation}
The first uses the comparator recurrence.  The second uses
\eqref{eq:hybrid-saturated-shift} in its successor step.  The
substitution $t=m=n$ is made into that already proved universally
parameterized equation; no parameter of an active induction hypothesis
is changed.  Consequently
\begin{equation}
 W(Sn,n)=B_n(n)=C(L(Sn,n),n,Sn)=n.
 \label{eq:hybrid-extra-step}
\end{equation}

Finally define the finite scalar term
\begin{equation}
 P(m)=I\bigl(m,0;u.C(L(SSu,m),u,Su)\bigr).
 \label{eq:hybrid-predecessor-term}
\end{equation}
Its zero equation is immediate.  To justify its successor equation
without assuming a separate extensionality rule for iterator bodies,
put
\[
 V(t,m)=I\bigl(t,0;u.C(L(SSu,Sm),u,Su)\bigr).
\]
The comparator recurrence identifies this update pointwise with
$B_m(u)$.  Equation induction on $t$, with $m$ fixed, gives
$V(t,m)=W(t,m)$: their base values are $0$, and the successor step
uses the pointwise update equation and then the induction equation.
Now \eqref{eq:hybrid-extra-step} yields
\[
 \boxed{\qquad P(0)=0,\qquad P(Sm)=V(Sm,m)=W(Sm,m)=m.\qquad}
\]

\subsection{Dependency boundary}

All terms in this section are built from $0,S,I$ by finite substitution.
All inductions have a single equation as their invariant.  Numerical
cases and equality transport are the already reconstructed logical
resources.  Successor injectivity may now be obtained by applying $P$
to an assumed equation $Sx=Sy$.  Equality tests, reflection of
quantifier-free formulas, pairing, and the recovery of general
primitive recursion can then be developed in the established order.
They were not used in order to manufacture predecessor.

This argument is a written internal derivation in the stated theory.
The finite implementation checks do not constitute a proof-assistant
formalization of this entire metatheorem.

% END reference/scalar_predecessor.tex


% BEGIN reference/equality_reflection.tex
% Standalone section.  Requires only amsmath and amssymb in the parent document.
% No theorem environments, macros, or bibliography keys are assumed.
% All arithmetic claims in this section are proved here.

\section{From literal-equation induction to full intuitionistic PRA}
\label{sec:eic-completion}\label{app:arithmetic}

This section supplies an unconditional completion theorem. It is adapted
with attribution from the current forward-left draft~\cite{forward}. In the
scalar application used here, every function occurring in the proof is
already a scalar term: predecessor is supplied by
Appendix~\ref{sec:hybrid-bootstrap}, the selector is
$C(n,a,b)=I(n,a;u.b)$, addition is $x+y=I(y,x;u.Su)$, and
truncated subtraction is $d(x,y)=I(y,x;u.Pu)$. All other functions below
are finite compositions of these. Thus this proof does not presuppose
that arbitrary primitive recursion has already been compiled. The QFI
conclusion first applies to the scalar language; the general PR compiler
comes afterward.  Literal-equation
induction, together with the primitive-recursive defining equations,
predecessor coverage, and separation of $0$ and $1$, proves induction for every
quantifier-free formula.  In particular, the proof does not assume a faithful
numerical equality test, stability or decidability of equality, successor
uniqueness, or any commutation law for an iterator.  We construct the equality
test and prove its reflection property for arbitrary elements of the
numerical domain.

\subsection{The exact hypotheses}

Work in ordinary intuitionistic first-order logic with equality and genuine
falsity $\bot$.  The language has one numerical sort and a symbol for each
fixed, finite primitive-recursive description.  Distinct descriptions are
not identified merely because they compute the same standard function.
Let $\mathsf{PRE}$ contain the defining equations for zero, successor,
projections, composition, and primitive recursion.  In particular, choose
descriptions $P$ and $C$ with equations
\begin{align}
 P(0)&=0, & P(Sx)&=x, \label{eq:eic-predecessor}\\
 C(0,a,b)&=a, & C(Sx,a,b)&=b. \label{eq:eic-selector}
\end{align}
Write $1=S0$.  In addition to $\mathsf{PRE}$, assume
\begin{align}
 \mathsf{Cov}:&\qquad
       \forall x\,\bigl(x=0\lor x=S(Px)\bigr),
       \label{eq:eic-coverage}\\
 \mathsf{Sep}:&\qquad 0=1\longrightarrow\bot,
       \label{eq:eic-separation}
\end{align}
and induction for each \emph{literal equation} $E(n,\vec p)$:
\begin{equation}
 \forall\vec p\,
 \left[
   \left(E(0,\vec p)\land
       \forall n\bigl(E(n,\vec p)\longrightarrow E(Sn,\vec p)\bigr)
   \right)
   \longrightarrow\forall n\,E(n,\vec p)
 \right].
 \label{eq:eic-induction}
\end{equation}
Here the schematic expression $E$ is a single equality of numerical
terms; the displayed universal quantifiers range only over numbers.
Every description and every instance of the schema is finite external
syntax.

The axiom formulation \eqref{eq:eic-induction} and the usual contextual
induction rule are interchangeable.  For example, if
\[
 \Gamma\vdash E(0,\vec p),
 \qquad
 \Gamma,E(z,\vec p)\vdash E(Sz,\vec p),
\]
where $z$ is absent from $\Gamma$ and the passive parameters, discharge the
second premise's equation and universally introduce $z$.  The axiom then
gives $\Gamma\vdash\forall n E(n,\vec p)$, which can be instantiated at any
target term.  Conversely, assume the antecedent of the axiom and apply the
contextual rule with fresh variables.  This uses the ordinary intuitionistic
logical rules and permits arbitrary, including quantified, formulas in
$\Gamma$.

\subsection{Constructor separation and two elementary tests}

Define the sign description and the countdown description by
\begin{equation}
 \operatorname{sg}(u)=C(u,0,1),
 \qquad
 K(u)=C(u,1,Pu).
 \label{eq:eic-sign-countdown}
\end{equation}
Their defining equations imply
\begin{align}
 \operatorname{sg}(0)&=0,
 &\operatorname{sg}(Sx)&=1,\label{eq:eic-sign-computation}\\
 K(0)&=1,
 &K(Sx)&=x.\label{eq:eic-countdown-computation}
\end{align}

\paragraph{Constructor laws.}
Successor is injective: applying $P$ to $Sx=Sy$ gives $x=y$ by
\eqref{eq:eic-predecessor}.  Successor never equals zero.  Indeed,
$Sx=0$ gives $1=0$ by applying $\operatorname{sg}$ and using
\eqref{eq:eic-sign-computation}; symmetry and
\eqref{eq:eic-separation} give $\bot$.  Consequently $1\ne0$, where
$A\ne B$ abbreviates $A=B\to\bot$.  None of these arguments uses
induction.

\paragraph{The sign test.}
For every $u$,
\begin{align}
 &\operatorname{sg}(u)=0\lor\operatorname{sg}(u)=1,
       \label{eq:eic-sign-range}\\
 &\operatorname{sg}(u)=0\quad\longleftrightarrow\quad u=0,
       \label{eq:eic-sign-zero}\\
 &\operatorname{sg}(u)=1\quad\longleftrightarrow\quad u\ne0.
       \label{eq:eic-sign-one}
\end{align}
To prove the range assertion, split \eqref{eq:eic-coverage}.  The branch
$u=0$ gives sign $0$, and the branch $u=S(Pu)$ gives sign $1$.
The implication $u=0\to\operatorname{sg}(u)=0$ is computation.
For its converse, the successor branch of coverage would make
$\operatorname{sg}(u)=0$ imply $1=0$, a contradiction; in that branch
explosion gives $u=0$ as well.  Similarly, if
$\operatorname{sg}(u)=1$ and $u=0$, computation gives $0=1$.
Conversely, $u\ne0$ excludes the zero branch of coverage, and the
successor branch computes the sign as $1$.

Coverage and \eqref{eq:eic-selector} also give, for all $u,v$,
\begin{equation}
 C(u,v,v)=v.
 \label{eq:eic-selector-idempotence}
\end{equation}
Both coverage branches compute to $v$, so this identity again needs no
induction.

\paragraph{The countdown test and the notation $W$.}
For every $u$,
\begin{equation}
 K(u)=0\quad\longleftrightarrow\quad u=1.
 \label{eq:eic-countdown-zero}
\end{equation}
If $u=1=S0$, \eqref{eq:eic-countdown-computation} gives $K(u)=0$.
Conversely, assume $K(u)=0$ and split coverage.  If $u=0$, then
$K(u)=1$, contradicting separation, and explosion gives $u=1$.
If $u=S(Pu)$, computation gives $K(u)=Pu$; hence $Pu=0$ and
$u=S0=1$.

If a presentation uses a description named $W$ with defining equations
\[
 W(0)=1,\qquad W(Sx)=x,
\]
then it proves $K(u)=W(u)$ for every $u$ already from coverage.
In the zero branch both sides are $1$; in the branch $u=S(Pu)$ both
sides are $Pu$.  Thus $K$ and $W$ are provably equal on the whole domain,
even if they are distinct descriptions.  This equality does not appeal to
standard-model agreement or to uniqueness of an arbitrary recurrence.

\subsection{Four elementary uses of equation induction}

Choose descriptions for addition and truncated subtraction, with recursion
on the second argument:
\begin{align}
 x+0&=x, &x+Sy&=S(x+y),\label{eq:eic-addition}\\
 d(x,0)&=x, &d(x,Sy)&=P(d(x,y)).\label{eq:eic-subtraction}
\end{align}
Using these equations and \eqref{eq:eic-induction}, we now derive
\begin{align}
 0+y&=y,\label{eq:eic-left-zero}\\
 (Sx)+y&=S(x+y),\label{eq:eic-left-successor}\\
 d(Sx,Sy)&=d(x,y),\label{eq:eic-cancel-successors}\\
 d(Sa,a)&=1.\label{eq:eic-diagonal}
\end{align}
Each induction below has a single equation as its induction formula.

For \eqref{eq:eic-left-zero}, induct on $y$.  The base is $0+0=0$.
The step is $0+Sy=S(0+y)=Sy$.

For \eqref{eq:eic-left-successor}, hold $x$ fixed and induct on $y$.
At zero both sides are $Sx$.  At a successor the induction hypothesis
gives
\[
 (Sx)+Sy=S\bigl((Sx)+y\bigr)
       =S\bigl(S(x+y)\bigr)=S(x+Sy).
\]
Combining \eqref{eq:eic-left-successor} with the defining equation for
addition gives the useful rearrangement
\begin{equation}
 (Sy)+r=y+Sr.
 \label{eq:eic-shift-successor}
\end{equation}
Indeed, both sides are $S(y+r)$.

For \eqref{eq:eic-cancel-successors}, hold $x$ fixed and induct on $y$.
The base computes as
\[
 d(Sx,S0)=P(d(Sx,0))=P(Sx)=x=d(x,0).
\]
Applying $P$ to the induction equation gives the step:
\[
 d(Sx,S(Sy))=P(d(Sx,Sy))=P(d(x,y))=d(x,Sy).
\]

Finally, induct on $a$ for \eqref{eq:eic-diagonal}.  The base is
$d(S0,0)=S0=1$.  At a successor,
\[
 d(S(Sa),Sa)=d(Sa,a)=1,
\]
using \eqref{eq:eic-cancel-successors} and the induction hypothesis.

\subsection{An equation invariant that reflects actual equality}

The remaining arithmetic ingredient is the following equation, valid for
all $x,y$:
\begin{equation}
 C\bigl(\operatorname{sg}(d(x,y)),\ x,\ y+d(x,y)\bigr)=x.
 \label{eq:eic-residual-invariant}
\end{equation}
Its meaning is that a nonzero subtraction residual reconstructs the
original argument by addition.  The selector expresses this assertion as
one literal equation, so equation induction applies directly.

Hold $x$ fixed and induct on $y$.  At $y=0$, subtraction gives residual
$x$, and \eqref{eq:eic-left-zero} gives $0+x=x$.  The desired equation
is consequently $C(\operatorname{sg}(x),x,x)=x$, which is
\eqref{eq:eic-selector-idempotence}.

For the step, assume \eqref{eq:eic-residual-invariant} at $y$ and
abbreviate
\[
 r=d(x,y),\qquad r'=Pr=d(x,Sy).
\]
These are local term abbreviations, not additional function symbols.
Split coverage for $r'$.

If $r'=0$, then $\operatorname{sg}(r')=0$, and the successor instance
of \eqref{eq:eic-residual-invariant} reduces immediately to $x=x$.

In the other branch, $r'=S(Pr')$, so constructor separation gives
$r'\ne0$.  Split coverage for $r$.  Its zero branch would imply
$r'=P0=0$, a contradiction, from which the desired step equation follows.
In its successor branch,
\[
 r=S(Pr)=Sr',\qquad \operatorname{sg}(r)=1.
\]
The induction equation therefore computes to $y+r=x$.  Meanwhile,
\eqref{eq:eic-shift-successor} gives
\[
 (Sy)+r'=y+Sr'=y+r=x.
\]
Since $r'$ is a successor, $\operatorname{sg}(r')=1$, so this is
exactly the branch selected in the successor instance of
\eqref{eq:eic-residual-invariant}.  This completes the equation induction.
All case distinctions came from coverage; no induction on an implication
or a disjunction was used.

\paragraph{Reflection at residual $1$.}
We have, for arbitrary $a,b$,
\begin{equation}
 d(Sa,b)=1\quad\longleftrightarrow\quad a=b.
 \label{eq:eic-residual-reflection}
\end{equation}
For the forward direction, substitute $x=Sa$, $y=b$ into
\eqref{eq:eic-residual-invariant}.  If $d(Sa,b)=1$, the sign is $1$
and the selector returns its second branch, giving
\[
 b+1=Sa.
\]
The defining equation for addition gives $b+1=S(b+0)=Sb$.  Thus
$Sb=Sa$, and applying $P$ gives $b=a$, hence $a=b$.
Conversely, $a=b$ and \eqref{eq:eic-diagonal} give
$d(Sa,b)=d(Sa,a)=1$.

Define the equality characteristic by the fixed description
\begin{equation}
 e(a,b)=\operatorname{sg}\bigl(K(d(Sa,b))\bigr).
 \label{eq:eic-equality-characteristic}
\end{equation}
The sign, countdown, and residual reflection lemmas give
\begin{align}
 e(a,b)=0&\quad\longleftrightarrow\quad a=b,
       \label{eq:eic-equality-reflection}\\
 e(a,b)=0&\ \lor\ e(a,b)=1.
       \label{eq:eic-equality-range}
\end{align}
In detail, $e(a,b)=0$ is equivalent to $K(d(Sa,b))=0$ by
\eqref{eq:eic-sign-zero}, which is equivalent to $d(Sa,b)=1$ by
\eqref{eq:eic-countdown-zero}, and then to $a=b$ by
\eqref{eq:eic-residual-reflection}.  The range assertion is
\eqref{eq:eic-sign-range} applied to $K(d(Sa,b))$.

The reflection property concerns the logic's actual equality throughout
the domain.  It was not assumed as part of a numerical comparison symbol,
and its proof did not restrict $a$ or $b$ to standard numerals.
It also yields decidable equality.  In the $e(a,b)=0$ branch,
\eqref{eq:eic-equality-reflection} gives $a=b$.  In the $e(a,b)=1$
branch, an assumption $a=b$ would give $e(a,b)=0$ and therefore $0=1$,
contradicting separation.  Hence
\begin{equation}
 \forall a\,\forall b\,\bigl(a=b\lor(a=b\to\bot)\bigr)
 \label{eq:eic-equality-decidable}
\end{equation}
is a theorem, not an extra hypothesis.

\subsection{Compilation of all quantifier-free formulas}

For each fixed quantifier-free formula $A(\vec p)$, construct a term
$b_A(\vec p)$ by external recursion on its finite syntax.  Take equality,
$\bot$, $\land$, $\lor$, and $\to$ as the primitive formula constructors;
negation is implication into $\bot$.  Set
\begin{align}
 b_{s=t}&=e(s,t), & b_{\bot}&=1,\label{eq:eic-boolean-atoms}\\
 b_{A\land D}&=C(b_A,b_D,1),\label{eq:eic-boolean-and}\\
 b_{A\lor D}&=C(b_A,0,b_D),\label{eq:eic-boolean-or}\\
 b_{A\to D}&=C(b_A,b_D,0).\label{eq:eic-boolean-implies}
\end{align}
Every term here is a finite composition of fixed primitive-recursive
descriptions; if desired, projections and composition give it a single
description symbol.  No internal quantification over formulas or programs
is involved.

An external induction on the syntax of $A$ proves, within the arithmetic
theory,
\begin{equation}
 \bigl(b_A=0\lor b_A=1\bigr)
 \quad\text{and}\quad
 \bigl(b_A=0\longleftrightarrow A\bigr).
 \label{eq:eic-boolean-correctness}
\end{equation}
Here are the full connective checks.  In every case $b_A=1$ implies
$\neg A$: if $A$ held, the already established implication
$A\to b_A=0$ would contradict separation.

\paragraph{Atomic formulas and falsity.}
For an equation, both claims are
\eqref{eq:eic-equality-reflection}--\eqref{eq:eic-equality-range}.
For falsity the range assertion is $1=0\lor1=1$, with its right disjunct
proved by reflexivity.  Also $1=0\to\bot$ follows from separation and
symmetry, while $\bot\to1=0$ follows by explosion.  Thus $1=0$ is
equivalent to $\bot$, as required.

\paragraph{Conjunction.}
Split $b_A=0\lor b_A=1$.  In the zero branch, $A$ holds and
$b_{A\land D}=b_D$.  Its range follows from the range of $b_D$,
and its being zero is equivalent to $D$, hence to $A\land D$ under
the known $A$.  In the one branch, $\neg A$ holds and
$b_{A\land D}=1$.  Its range is immediate.  If it were zero,
separation would give $\bot$, hence $A\land D$; conversely an
assumption $A\land D$ contradicts $\neg A$, and explosion gives
$b_{A\land D}=0$.

\paragraph{Disjunction.}
In the branch $b_A=0$, the formula $A$ holds and
$b_{A\lor D}=0$.  The range is immediate, and both $b_{A\lor D}=0$
and $A\lor D$ hold, proving their equivalence.  In the branch
$b_A=1$, we have $\neg A$ and $b_{A\lor D}=b_D$.  Its range follows
from the range of $b_D$.  If its value is zero, then $D$ and hence
$A\lor D$.  Conversely, from $A\lor D$, the $A$ alternative
contradicts $\neg A$ and yields $b_D=0$ by explosion, while the
$D$ alternative yields $b_D=0$ directly.  This proves the required
equivalence in both alternatives.

\paragraph{Implication.}
In the branch $b_A=0$, the formula $A$ holds and
$b_{A\to D}=b_D$.  Its range follows from the range of $b_D$.
If its value is zero, then $D$, which gives $A\to D$.  Conversely,
$A\to D$ and the known $A$ give $D$, hence $b_D=0$.
In the branch $b_A=1$, we have $\neg A$ and
$b_{A\to D}=0$.  Its range is immediate.  To prove $A\to D$,
assume $A$; this contradicts $\neg A$, and explosion gives $D$.
Thus both $b_{A\to D}=0$ and $A\to D$ hold in this branch, and
their equivalence follows.  This argument uses the proved two-value
range of the finite characteristic.  It does not use excluded middle
for arbitrary formulas or change the surrounding intuitionistic logic.

This completes the external syntactic induction proving
\eqref{eq:eic-boolean-correctness} for every quantifier-free formula.

\subsection{Recovery of induction and the full theory}

Let $A(n,\vec p)$ be any fixed quantifier-free formula.  Suppose an
arbitrary retained context $\Gamma$ proves its base and its induction
step:
\[
 \Gamma\vdash A(0,\vec p),
 \qquad
 \Gamma,A(z,\vec p)\vdash A(Sz,\vec p),
\]
with the usual freshness condition on $z$.  By
\eqref{eq:eic-boolean-correctness}, the base gives
$\Gamma\vdash b_A(0,\vec p)=0$.  Under the additional assumption
$b_A(z,\vec p)=0$, the same equivalence gives $A(z,\vec p)$,
the given step gives $A(Sz,\vec p)$, and the equivalence gives
$b_A(Sz,\vec p)=0$.  Apply literal-equation induction to
$b_A(n,\vec p)=0$ and then apply its reflection implication.
The result is $\Gamma\vdash A(t,\vec p)$ for every target term $t$.
Thus the theory proves the full contextual quantifier-free induction
rule, including in arbitrary quantified contexts.

\paragraph{Completion theorem.}
In the common language of fixed primitive-recursive descriptions, the
theory
\[
 \mathsf{PRE}+\mathsf{Cov}+\mathsf{Sep}+\mathsf{EI}
\]
has exactly the same sequents as ordinary first-order intuitionistic
primitive recursive arithmetic with all quantifier-free induction
instances.  The forward simulation follows from the recovered
quantifier-free induction, successor injectivity, and successor separation;
all primitive-recursive defining equations and all logical rules are
shared.  For the reverse simulation, literal-equation induction is an
instance of quantifier-free induction, and successor separation gives
$\mathsf{Sep}$.  Quantifier-free induction proves $\mathsf{Cov}$:
at zero choose its left disjunct, and at a successor use
$Sx=S(P(Sx))$, which follows from $P(Sx)=x$, for its right disjunct.

Both interpretations therefore use the full domain, defined by $x=x$,
and preserve every function symbol, equality, logical connective, and
quantifier.  Both composites have the identity isomorphism, with graph
$x=y$.  The theories are consequently identity-domain bi-interpretable
and equiconsistent.  This completion theorem is unconditional once the
literal-equation induction rule has been obtained; it has no dependency
on a classical successor-uniqueness reduction.

% END reference/equality_reflection.tex


% BEGIN reference/primitive_recursion.tex
\section{Pairing, primitive recursion, and internal numerical composites}
\label{sec:iterator-coding}
This appendix is adapted with attribution from the progress-closure
arithmetic development~\cite{progress}. Its prerequisites have now been
proved: scalar predecessor, arithmetic reflection, and full quantifier-free
induction. Pairing is a derived proof tool; it is not part of any checker.

\subsection{Pairing from a bounded square-root search}

Once QFI is available, a familiar square-based injection gives a
particularly direct compiler. Write $a\le b$ for $d(a,b)=0$.
The usual linear-order and semiring identities follow by fixed
quantifier-free inductions from the successor, addition, and monus
equations. Multiplication itself is the scalar term
\[
                        a\cdot b=I(b,0;u.u+a).
\]
Only primitive-recursive arithmetic is involved in the following
bounded root search; there is no minimization over an unbounded domain.

For a fixed parameter $n$, put
\[
 U_n(v)=C(d((Sv)^2,n),Sv,v),\qquad
 R_n(i)=I(i,0;v.U_n(v)),\qquad q(n)=R_n(n).
\]
The update increments $v$ exactly when $(v+1)^2\le n$, and
otherwise retains it. QFI in $i$, with $n$ fixed, proves
\[
 R_n(i)\le i,\qquad R_n(i)^2\le n,\qquad
 R_n(i)=i\ \vee\ n<(S(R_n(i)))^2.              \tag{Root invariant}
\]
The base has $R_n(0)=0$. If the update increments, its test
contradicts the strict-upper-bound alternative in the third invariant.
Thus the old state equals $i$, and the new state equals $Si$;
its square bound is the increment test. If the update does not increment, the
strict upper square bound holds and the state remains fixed.
An already established strict upper square bound forces every
subsequent update to remain fixed. These are exactly the cases
needed for the three displayed invariants.

At $i=n$, either the strict upper square bound already holds or
$R_n(n)=n$, in which case $n<(n+1)^2$ supplies it. Thus
\[
                        q(n)^2\le n<(q(n)+1)^2.           \tag{Root}
\]
Monotonicity of squares shows uniqueness of a number satisfying
Root. All formulas just used are quantifier free in the scalar
language, with fixed $n$ during the induction.

Now define
\[
 \pi(a,b)=(a+b)^2+a,\qquad
 L(n)=d(n,q(n)^2),\qquad H(n)=d(q(n),L(n)).
\]
For $k=a+b$, one has $k^2\le k^2+a<(k+1)^2$ since $a\le k$.
Root uniqueness gives $q(\pi(a,b))=a+b$, and therefore
\[
                    L(\pi(a,b))=a,\qquad H(\pi(a,b))=b.
\]
Surjectivity of $\pi$ is unnecessary. The original numerical domain
is never replaced by the range of this coding; pairing is used only
to store the finite state of a fixed primitive recursion. This
standard square pairing is a convenient reconstruction here, not
a claimed new arithmetic basis.

\subsection{Compiling each primitive recursion}

Suppose the fixed target description is
\[
 f(0,\vec p)=g(\vec p),\qquad
 f(Sn,\vec p)=h(n,f(n,\vec p),\vec p).
\]
Assume by external structural induction that the earlier $g,h$
already have scalar translations $g^*,h^*$. Define
\[
 W(u,\vec p)=\pi(SL(u),h^*(L(u),H(u),\vec p)),
\]
\[
 V(n,\vec p)=I(n,\pi(0,g^*(\vec p));u.W(u,\vec p)),
 \qquad f^*(n,\vec p)=H(V(n,\vec p)).
\]
The update has only its state variable bound, with every $p_i$
fixed. Equation induction proves $L(V(n,\vec p))=n$.
The two projection equations and scalar computation then prove
both defining equations for $f^*$ internally. Initial functions
and composition translate directly. This completes the translation
of every fixed finite PR description without adding a universal
evaluator.

\subsection{Why both numerical round trips are internal}

Two fixed scalar iterators with the same seed and provably equal
updates are internally equal at all counters: use literal-equation
induction comparing their values. Their base agrees, and their
successor equations plus congruence prove the step. The same
uniqueness argument works for ordinary primitive recursion.

For an ordinary PR description, induct externally on its finite
construction. The double translation satisfies its original
defining equations, so internal equation induction identifies it
with the original operation. For a native scalar description,
earlier translated updates agree by the external induction
hypothesis; scalar uniqueness identifies the double-translated
iterator. These proofs hold at arbitrary internal counters.

Thus the domain relation $x=y$ is provably a total bijection
preserving every fixed numerical operation in both composites.
The formula constructions and round trips in
Section~\ref{sec:choice-introduction} give the corresponding formula
composites. Both interpretations preserve falsity, taking native
$0=1$ to an ordinary contradiction and ordinary $\bot$ to native
$0=1$. Finite proof translations therefore establish equiconsistency.

% END reference/primitive_recursion.tex


\clearpage
\begin{thebibliography}{99}
\raggedright
\bibitem{corpus}
Vladimir Reshetnikov, repository maintainer.
\emph{ProveIt, current draft collection}, commit
\file{fb2dbd2a33c6b132671a54b029d0bf3cd98c633d}.
\url{https://github.com/VladimirReshetnikov/ProveIt/tree/fb2dbd2a33c6b132671a54b029d0bf3cd98c633d/docs/drafts}.
\bibitem{arity}
\emph{An Arity Calculus for iPRA},
\file{ipra_arity_calculus_2026-11-10.zip} in \cite{corpus}.
Source of the unary/triplicate representation, arity opening, and the
five attributed mathematical appendices.
\bibitem{uniform}
\emph{Uniform Predecessor Certificates for Intuitionistic PRA},
\file{ipra_uniform_predecessor_2026-11-10.zip} in \cite{corpus}.
Source of predecessor packets, supplied frontier implementations, and
the named producer adapted here.
\bibitem{constructed}
\emph{Constructed Computation},
\file{ipra_constructed_computation_2026-11-10.zip} in \cite{corpus}.
\bibitem{generated}
\emph{Numerical Axioms as Local Assumptions}, inherited in the reference
material of \cite{arity,constructed}.
\bibitem{transient}
\emph{A One-State Recursor for iPRA},
\file{ipra_transient_seed_2026-11-10.zip} in \cite{corpus}.
\bibitem{affine}
\emph{Affine Depths for a Small iPRA Checker},
\file{ipra_affine_arity_2026-11-10.zip} in \cite{corpus}.
\bibitem{scopes}
\emph{Local Scopes}, \file{ipra_local_scopes_2026-11-10.zip} in
\cite{corpus}. Includes the earlier relational obstruction for a
tied-seed basis.
\bibitem{salehi}
Saeed Salehi. \emph{Decidable Formulas of Intuitionistic Primitive
Recursive Arithmetic}. Reports on Mathematical Logic \textbf{36}
(2002), 55--61.
\url{https://rml.tcs.uj.edu.pl/rml-36/36-SALEHI.pdf}.
\bibitem{erasure}
Inherited scalar predecessor development, retained with attribution in
\file{reference/scalar_predecessor.tex} of \cite{arity}.
\bibitem{forward}
Inherited forward-left equality-reflection development, retained in
\file{reference/equality_reflection.tex} of \cite{arity}.
\bibitem{progress}
Inherited progress-closure pairing and primitive-recursion compiler,
retained in \file{reference/primitive_recursion.tex} of \cite{arity}.
\end{thebibliography}
\end{document}
